\documentclass[11pt]{article}
\usepackage[T1]{fontenc}
\usepackage[utf8]{inputenc}
\usepackage{lmodern}
\usepackage{microtype}
\usepackage[letterpaper,margin=0.92in]{geometry}
\usepackage{amsmath,amssymb,amsthm,mathtools}
\usepackage{booktabs,longtable,tabularx,array,multirow}
\usepackage{enumitem}
\usepackage{xcolor}
\usepackage{xurl}
\usepackage{hyperref}
\usepackage[nameinlink,noabbrev]{cleveref}
\usepackage{listings}
\usepackage{fancyhdr}
\usepackage{graphicx}
\usepackage{enumitem}

\hypersetup{
  colorlinks=true,
  linkcolor=blue!45!black,
  citecolor=green!35!black,
  urlcolor=blue!55!black,
  pdftitle={A Certified Finitary Presentation Framework for Mathematics in iPRA},
  pdfauthor={OpenAI, prepared for Vladimir Reshetnikov}
}
\setlength{\emergencystretch}{3em}
\setlength{\headheight}{14pt}
\setlist{nosep,leftmargin=2em}
\pagestyle{fancy}
\fancyhf{}
\fancyhead[L]{Certified finitary presentations in iPRA}
\fancyhead[R]{\thepage}
\renewcommand{\headrulewidth}{0.4pt}

\newtheorem{theorem}{Theorem}[section]
\newtheorem{lemma}[theorem]{Lemma}
\newtheorem{proposition}[theorem]{Proposition}
\newtheorem{corollary}[theorem]{Corollary}
\newtheorem{definition}[theorem]{Definition}
\newtheorem{principle}[theorem]{Principle}
\newtheorem{warning}[theorem]{Warning}
\newtheorem{example}[theorem]{Example}
\newtheorem{remark}[theorem]{Remark}
\theoremstyle{definition}
\newtheorem{construction}[theorem]{Construction}

\newcommand{\N}{\mathbb N}
\newcommand{\Z}{\mathbb Z}
\newcommand{\Q}{\mathbb Q}
\newcommand{\R}{\mathbb R}
\newcommand{\C}{\mathbb C}
\newcommand{\ipra}{\mathrm{iPRA}}
\newcommand{\PRA}{\mathrm{PRA}}
\newcommand{\PR}{\mathrm{PR}}
\newcommand{\CFP}{\mathrm{CFP}}
\newcommand{\eps}{\varepsilon}
\newcommand{\bits}{\operatorname{bits}}
\newcommand{\htop}{\operatorname{ht}}
\newcommand{\dist}{\operatorname{dist}}
\newcommand{\code}[1]{\nolinkurl{#1}}
\newcommand{\abs}[1]{\left|#1\right|}
\newcommand{\norm}[1]{\left\lVert#1\right\rVert}
\newcommand{\ceil}[1]{\left\lceil#1\right\rceil}
\newcommand{\floor}[1]{\left\lfloor#1\right\rfloor}
\newcommand{\angles}[1]{\left\langle#1\right\rangle}
\newcommand{\repo}{\texttt{VladimirReshetnikov/ProveIt}}

\lstdefinestyle{leanstyle}{
  basicstyle=\ttfamily\small,
  columns=fullflexible,
  keepspaces=true,
  frame=single,
  breaklines=true,
  showstringspaces=false,
  keywordstyle=\color{blue!55!black},
  commentstyle=\color{green!35!black}
}
\lstdefinestyle{pythonstyle}{
  basicstyle=\ttfamily\small,
  columns=fullflexible,
  keepspaces=true,
  frame=single,
  breaklines=true,
  showstringspaces=false,
  keywordstyle=\color{blue!55!black},
  commentstyle=\color{green!35!black}
}

\title{\bfseries A Certified Finitary Presentation Framework\\
for Mathematics in Intuitionistic Primitive Recursive Arithmetic\\[0.5em]
\large A common architecture for finite algebra, represented analysis,\\
geometry, probability, number theory, and constructive obstruction audits}
\author{Prepared for Vladimir Reshetnikov}
\date{15 November 2026}

\begin{document}
\maketitle

\begin{abstract}
The proof packets under
\code{Logic/PrimitiveRecursiveArithmetic} exhibit a stable positive pattern
and a stable family of failures.  Finite exact algebra, bounded search, finite
sums, grid arguments, and explicit descents usually fit intuitionistic
primitive recursive arithmetic once their loop invariants are written as
quantifier-free Boolean equations.  Difficulties arise when a textbook proof
silently quantifies over arbitrary functions, invokes an unspecified limit or
compactness principle, chooses an exact branch at a degenerate input, hides an
ineffective constant, or stops at an ambient Lean theorem instead of producing
an object-language derivation of type \code{Prf PRAAx}.

This paper develops a unified \emph{Certified Finitary Presentation} (\(\CFP\))
framework.  A mathematical object is presented by finite observations,
primitive-recursive precision transformers, domain and apartness certificates,
and bounded computation traces.  The central finite-trace reflection theorem
reduces correctness to a base equation, a quantifier-free step invariant, and
a final checker implication.  Regular names, finite nets, nested isolating
certificates, robust overlapping branches, and height/covering certificates
are then reusable modules rather than theorem-specific prose conventions.

The framework yields fully explicit constructive replacements for several
reported obstacles: all-precision induction is replaced by a finite query
table at a free precision parameter; exact branch selection is replaced by an
overlapping decision or by a decision under a supplied gap; abstract
compactness is replaced by primitive-recursive nets; exact roots and axes are
represented by coherent isolating data; and Diophantine finiteness is factored
into a finite verifier plus the genuinely mathematical production of a height
or covering certificate.  It also proves sharp boundaries: there is no total
universal primitive-recursive evaluator; a bound on the number of solutions
does not provide a search cutoff; and exact selections known to imply LLPO,
LPO, or Markov principles cannot be recovered by a change of encoding.

A detailed Lean plan is given against the existing declarations
\code{PR}, \code{PTm}, \code{PRAAx}, \code{PRA.qfInd}, \code{NFm},
\code{NPrf}, and \code{pra_iff_native}.  The accompanying archive contains
LaTeX and PDF, a machine-readable difficulty matrix, a Lean semantic prototype
and object-language blueprint, and an exact Python certificate kernel with
passing regression tests.  The Lean files are an explicit completion plan, not
a claim that the generic embedded reflection theorem has already been
kernel-checked.
\end{abstract}

\tableofcontents
\newpage

\section{Scope, standard of proof, and snapshot}
\label{sec:scope}

\subsection{The target is an arithmetic derivation}

The phrase ``formalizable in \(\ipra\)'' is used here in a strict sense.  The
final mathematical assertion must be a first-order formula over natural-number
codes.  Its functions are fixed primitive-recursive symbols, its logic is
intuitionistic, and every use of induction has a quantifier-free induction
formula.  A completed Lean endpoint therefore has the shape
\[
   \Gamma\;\vdash[\texttt{PRAAx}]\;\varphi,
\]
not merely an ambient theorem about Lean's \texttt{Nat}, \texttt{Rat}, or
\texttt{Real}.  The latter can validate an algorithm or an algebraic identity,
but it does not by itself construct a derivation in the embedded theory.

The existing \code{ConciseIPRA} development supplies both sides of this
boundary.  Its ordinary presentation has the types \code{PR} and \code{PTm},
the axiom family \code{PRAAx}, and the derived rule \code{PRA.qfInd}.  Its
native scalar calculus has iterator terms, equation formulas, the two binders
\code{Q} and \code{D}, and derivations \code{NPrf}.  The theorems
\code{pra_iff_native} and \code{native_iff_pra} faithfully transport an
already completed arithmetic derivation.  They are bridges, not automatic
mathematics compilers.

\subsection{Repository snapshot and audit scope}

The public repository was inspected at main-branch commit
\begin{center}
\small\ttfamily
7d492fbfc5fee3b2b9dec0b2623b0b82effe05f2.
\end{center}
The proof-packet index at that snapshot records forty-eight delivered packets
in algebra, analysis, analytic and Diophantine number theory, combinatorics,
geometry, probability, and topology.  The index also warns that packets were
filed as delivered: placement did not rerun their code or compile their Lean
files.  Some contain complete finite arguments; some contain uncompiled
blueprints, explicit \code{sorry} obligations, uninstantiated analytic
contracts, or deliberately isolated impossibility results.

This paper is therefore not a blanket certification of the packets.  It uses
them as a large empirical sample from which to extract common proof shapes and
common failure modes.  Statements about a particular packet are limited to the
paths listed in \cref{app:sources}.  The positive theorems proved below are
framework theorems: finite-trace reflection, regular-name calculus, robust
branching, finite-net search, and certificate transformations.  They do not
supply the missing deep analytic or arithmetic estimates in a packet whose
report explicitly leaves them open.

\subsection{Three notions that must not be conflated}

There are three useful but distinct levels of completion.

\begin{definition}[Semantic implementation]
A semantic implementation is executable code or an ambient Lean theorem that
computes the intended finite object and proves a proposition about its ordinary
meaning.
\end{definition}

\begin{definition}[Arithmetic proof architecture]
An arithmetic proof architecture specifies natural-number codes, fixed
primitive-recursive functions, exact formulas, induction invariants, bounds,
and a dependency-ordered route to an embedded derivation.
\end{definition}

\begin{definition}[Embedded completion]
An embedded completion is an actual derivation in \code{PRAAx}, optionally
transported to \code{NPrf}, with an axiom audit of the final Lean declaration.
\end{definition}

A packet may be strong at the first two levels while incomplete at the third.
The framework is designed to make the transition from the second to the third
mechanical enough that the remaining gap is visible and localized.

\subsection{What ``general'' can mean in a first-order theory}

A textbook theorem often begins ``for every continuous function'' or ``for
every real random variable.''  Literal higher-type quantification is not
available in first-order \(\ipra\).  There are three legitimate replacements.

\begin{enumerate}
\item A \emph{fixed-evaluator schema}: choose a primitive-recursive evaluator
      \(a(\mathbf p,k)\) externally, add its regularity or domain assumptions,
      and prove the resulting first-order formula uniformly in the numerical
      parameters \(\mathbf p\).
\item A \emph{finite-data theorem}: make a rational table, finite law,
      polynomial, graph, triangulation, or certificate the numerical input.
      This is a single uniform first-order theorem.
\item A \emph{trace-relative theorem}: input a finite computation or proof
      trace and prove that acceptance by a primitive-recursive verifier implies
      the desired arithmetic conclusion.  Trace production may be specialized,
      while verification remains uniform.
\end{enumerate}

The last form is central.  It gives generic infrastructure without postulating
a total universal evaluator for all primitive-recursive programs, which is
impossible by diagonalization (\cref{thm:no-universal-pr}).

\subsection{Contribution and limits}

The main contribution is a reusable normal form for positive proofs:
\[
\boxed{\begin{gathered}
\text{fixed numerical input and precision}\;\longrightarrow\;\text{bounded trace}\\
\longrightarrow\;\text{Boolean invariant}\;\longrightarrow\;\text{checked observation}
\end{gathered}}
\]
The exact or asymptotic theorem is recovered only after the finite observation
has been proved for an arbitrary free precision parameter.  This order is the
logical discipline that prevents a hidden use of completeness, countable
choice, or induction on a quantified Cauchy assertion.

The framework resolves many \emph{formulation and compilation} difficulties.
It does not turn ineffective constants into effective ones, produce rational
representatives from abstract Selmer data, or uniformly choose an axis at the
identity rotation.  In those cases it identifies a minimal positive theorem
and isolates the missing producer or the constructive obstruction.

\section{Diagnosis of the existing proof packets}
\label{sec:diagnosis}

\subsection{The successful pattern}

Across the packets, the arguments most naturally internalized in \(\ipra\)
share six properties.

\begin{enumerate}
\item Every object used by the calculation has a finite code: an integer,
      rational, finite list, polynomial, matrix, finite probability law,
      triangulation, or bounded computation state.
\item Every loop has a displayed primitive-recursive bound.
\item Its invariant is an equality or inequality of coded finite objects.
      Bounded universal assertions are compressed into Boolean folds.
\item Every analytic symbol is backed by a finite approximant and an explicit
      error estimate; no appeal is made to an abstract completion axiom.
\item Every division, normalization, pivot, or chart carries a positive
      denominator or apartness certificate.
\item The desired theorem is obtained by introducing the precision variable
      only after the finite calculation has been proved.
\end{enumerate}

Finite polynomial identities are the model case.  Normalize both sides to a
coefficient list, verify coefficient equality, and prove evaluation respects
normalization by induction on the finite list.  The same structure appears in
finite sums, product telescoping, exact rational probability, grid parity,
bounded point searches, and descent traces.

\subsection{The recurrent failure modes}

The negative pattern is equally stable.  The following table gives the
framework diagnosis, not a verdict on the mathematical truth of the classical
statement.

\begin{longtable}{@{}p{0.20\textwidth}p{0.33\textwidth}p{0.39\textwidth}@{}}
\toprule
Failure mode & Symptom & Correct repair or boundary \\
\midrule
\endhead
Higher-type quantification & ``For every real/function/measure'' is written as
if the function itself were a first-order variable. & Use a fixed evaluator
schema, finite input code, or finite trace relation.  Never postulate a total
universal PR evaluator.\\[3pt]
All-precision induction & The induction predicate says that a sequence is
Cauchy, equal, or correct at every precision. & Fix the output precision;
collect the finitely queried samples in a Boolean \(\mathsf{GoodTable}\);
induct on a quantifier-free finite invariant; quantify precision last.\\[3pt]
Exact branch at a boundary & A proof chooses a sign, pivot, eigenvector, root
ordering, argument, or chart when the datum may lie exactly on the boundary. &
Use overlapping branches, homogeneous/projective output, or a supplied
positive gap.  If exact selection implies LLPO/LPO/Markov, retain the
obstruction.\\[3pt]
Unqualified partial operation & Division, square root, logarithm, or
normalization is treated as total without carrying its domain proof. & Define a
totalized code-level operation and prove the intended law under a Boolean
domain/apartness guard.\\[3pt]
Abstract compactness & A convergent subsequence, minimizer, root, or finite
subcover is selected without a modulus or net. & Supply a primitive-recursive
finite net and continuity modulus; return approximate witnesses or coherent
nested certificates.\\[3pt]
Ineffective threshold & A classical proof says ``there exists a constant'' but
no computable bound precedes later parameters. & Factor the result into a
finite verifier plus an explicit analytic-input certificate.  Proving a
producer is separate mathematics.\\[3pt]
Cardinality in place of location & Finiteness or a point-count bound is used as
if it supplied a search cutoff. & Require a height bound, complete list, or
covering representatives.  Cardinality alone has no stopping rule.\\[3pt]
Ambient/formal gap & Lean proves an ordinary semantic theorem, but no
\code{Prf PRAAx} object is constructed. & Implement the fixed PR symbols,
Boolean checker, and QF invariant; then use the existing faithful translation.\\
\bottomrule
\end{longtable}

\subsection{Representative packets}

The diagnosis is especially visible in the following reports.

\paragraph{Euler--Rodrigues.}
The unrestricted exact unit-axis selector for arbitrary represented rotations
is reported to imply LLPO.  The same packet obtains rational exact axes,
separated exact axes, arbitrary-precision approximate axes, and uniform
quaternion data.  This is the canonical example of replacing a discontinuous
selector by robust data rather than weakening arithmetic by an extra logical
principle.

\paragraph{Dirichlet and Kronecker approximation.}
Sharp closed endpoints, a real-valued denominator bound, and density from
negative irrationality alone lead respectively to LLPO, LPO, and a Markov-type
principle.  Strict inequalities with margins and positive separation moduli
survive.  Thus ``add a quantitative margin'' is not cosmetic; it changes a
non-realizable selection problem into finite rational comparison.

\paragraph{Faltings, Siegel, and Mordell--Weil.}
These reports isolate missing height bounds or covering representatives.
Finite search below a supplied height is straightforward; producing the height
is the hard theorem.  A complete point list and a height bound are mutually
convertible in the coded setting, while a mere cardinality bound is not.
Mordell--Weil descent becomes constructive relative to actual quotient
representatives, but finite abstract quotient data do not automatically
produce rational points.

\paragraph{Prime distribution.}
Dirichlet and finite sieve inequalities can be made explicit.  The full
Siegel--Walfisz and Bombieri--Vinogradov routes retain ineffective or
exceptional-zero input.  The correct modular theorem checks a supplied
zero-free/exceptional certificate and propagates its constants; it must not
pretend that verification also produces the certificate.

\paragraph{Central limit theorems.}
Finite rational laws and finite product spaces give exact sums and exact
independence.  Infinite or continuous laws require tail, quantization, and
continuity moduli.  Rational thresholds avoid undecidable comparison with an
irrational normalized sum; transfer to a CDF metric uses a continuity or
anti-concentration estimate.

\subsection{A unifying interpretation}

These examples suggest that the fundamental object is not a ``real number''
or ``continuous map'' in isolation.  It is a bundle of finite interfaces:

\[
\begin{gathered}
\text{code and domain guard}\quad\text{finite observations}\quad
\text{error or residual}\\
\text{modulus}\quad\text{bounded trace}\quad
\text{soundness certificate}
\end{gathered}
\]

The \(\CFP\) framework makes this bundle explicit.  Once it is explicit, the
same arithmetic lemmas support analysis, geometry, probability, and number
theory, while the genuinely nonconstructive or ineffective step is no longer
hidden inside notation.

\section{The arithmetic base and its usable normal forms}
\label{sec:base}

\subsection{Ordinary iPRA}

We use the repository's ordinary first-order presentation.  For each arity
\(r\), \code{PR r} contains symbols generated from zero, successor,
projections, composition, and primitive recursion.  Terms \code{PTm n} are
built from \(n\) variables and these symbols.  The theory \code{PRAAx}
contains the defining equations, successor nonzeroness and injectivity, and
induction for quantifier-free formulas.  Intuitionistic natural deduction is
provided by \code{Deduction.lean}.

A bounded algorithm belongs to the theory in two different senses.  First, its
output function is a fixed \(\PR\) symbol because bounded loops and finite
folds are primitive recursive.  Second, its correctness must be proved.  That
proof is usually an instance of \code{PRA.qfInd}; primitive recursiveness alone
never implies correctness or truth.

\subsection{The native scalar core}

The native terms are generated by variables, \(0\), successor, and a scalar
iterator
\[
 I(n,a;u.h),\qquad
 I(0,a;u.h)=a,\qquad
 I(Sn,a;u.h)=h[I(n,a;u.h)/u].
\]
Native formulas are equations and two binders:
\[
 Q_y B(y)\equiv B(0)\to\forall y\,B(y),\qquad
 D_{x,y}B(x,y)\equiv\exists x\,(B(x,0)\to\forall y\,B(x,y)).
\]
The existing development derives ordinary intuitionistic logic and proves
faithful round trips between the native calculus, an equality theory, and
\(\ipra\).  The certificate layer proposed below can therefore target ordinary
\code{PRAAx}, where coding is readable, and export the finished theorem to the
smaller native calculus at the end.

\subsection{Booleanization of bounded assertions}

Let \(p(\mathbf a,i)\in\{0,1\}\) be a primitive-recursive characteristic
function.  Define
\[
\begin{aligned}
 \mathsf{All}_p(\mathbf a,0)&=1,\\
 \mathsf{All}_p(\mathbf a,n+1)&=
   \mathsf{All}_p(\mathbf a,n)\cdot p(\mathbf a,n),\\
 \mathsf{Any}_p(\mathbf a,0)&=0,\\
 \mathsf{Any}_p(\mathbf a,n+1)&=
   \max(\mathsf{Any}_p(\mathbf a,n),p(\mathbf a,n)).
\end{aligned}
\]
The intended bounded quantifiers become equations
\[
 \forall i<n\,P(i)\quad\rightsquigarrow\quad
 \mathsf{All}_p(\mathbf a,n)=1,
 \qquad
 \exists i<n\,P(i)\quad\rightsquigarrow\quad
 \mathsf{Any}_p(\mathbf a,n)=1.
\]

\begin{lemma}[Boolean-fold correctness]
\label{lem:boolean-fold}
For a fixed characteristic symbol \(p\), \(\ipra\) proves the usual
introduction and elimination rules for \(\mathsf{All}_p\) and
\(\mathsf{Any}_p\), uniformly in all numerical parameters.
\end{lemma}

\begin{proof}
For \(\mathsf{All}\), induction on \(n\) uses the quantifier-free invariant
that the accumulator equals one exactly when all processed bits equal one.
At the successor step, a product of two Boolean values is one exactly when
both are one.  Elimination at a particular \(i<n\) may be proved either by a
second bounded scan or by an invariant storing the first failure index.
For \(\mathsf{Any}\), use maximum or Boolean disjunction and store the least
successful index when witness extraction is required.  All state components
are natural codes and every formula in the induction is an equation between
primitive-recursive terms.  Hence \code{PRA.qfInd} applies.
\end{proof}

This elementary lemma is the main syntactic device for keeping invariants
quantifier free.  A phrase such as ``every previously processed coefficient
is correct'' should almost always mean an equation asserting that a finite
Boolean fold is one.

\subsection{Totalization and guarded laws}

The natural-number universe has no partial functions.  Every partial
mathematical operation is implemented by a total function plus a domain
predicate.  For example, rational division may return zero when the divisor
code is invalid or zero, while its theorem states
\[
 \mathsf{Rat}(x)=1\wedge \mathsf{Rat}(y)=1\wedge
 \mathsf{Nonzero}(y)=1
 \;\longrightarrow\;
 x/y\text{ has the expected cross-multiplication law}.
\]
This is preferable to trying to make every term intrinsically well typed.
The domain bits become fields of \(\mathsf{GoodTable}\), and invalid behavior
is irrelevant under the theorem's antecedent.

\subsection{A relation may be universal when a function cannot be}

A total primitive-recursive evaluator \(U(e,x)\) enumerating all unary
primitive-recursive functions cannot exist.  Nevertheless, a finite relation
\(\mathsf{EvalTrace}(e,x,t,y)=1\) can check that \(t\) is a well-formed
calculation tree for program code \(e\) on input \(x\), ending in \(y\).
Checking a finite tree is primitive recursive.  What fails to be primitive
recursive is a uniform function that always produces a terminating trace for
every code.

This relation/function distinction allows a genuinely general certificate
language.  A theorem may be uniform in \(e,x,t,y\) and prove
\[
 \mathsf{EvalTrace}(e,x,t,y)=1\to\mathsf{Post}(e,x,y)=1.
\]
For a fixed program, a separate compiler or proof constructs the trace or
replaces the relation by a direct \(\PR\) symbol.  No diagonal contradiction
arises because the verifier does not claim total trace production.

\section{Certified Finitary Presentations}
\label{sec:framework}

\subsection{Definition of a presentation}

A \emph{Certified Finitary Presentation} is a proof-oriented interface rather
than a new foundational ontology.

\begin{definition}[CFP object]
A \(\CFP\) object family with numerical parameter \(p\) consists of:
\begin{enumerate}
\item a decidable code-domain predicate \(\mathsf{Dom}(p)=1\);
\item finite observations \(\mathsf{obs}(p,k)\), where \(k\) is a precision,
      scale, truncation, or grid parameter;
\item a rational or integer error radius \(\mathsf{err}(p,k)\);
\item a regularity, consistency, or residual contract relating observations;
\item primitive-recursive bounds on all finite tables and traces used to
      compute or verify an observation.
\end{enumerate}
The evaluator is fixed at the meta-level, or its individual computation is
supplied by a finite trace.
\end{definition}

\begin{definition}[Certified transducer]
A certified transducer from presentations \(X\) to \(Y\) consists of a total
primitive-recursive output algorithm, a primitive-recursive precision
transformer \(\mu\), a bounded trace system, and a soundness theorem saying
that valid input observations at the finitely many precisions requested by
\(\mu\) imply the requested output observation and error bound.
\end{definition}

The transducer need not consume a single sample.  It may request an adaptive
finite table: later indices can depend on earlier rational comparisons.  The
complete run is still a bounded trace.  Its input contract records every query
actually made, and its soundness proof is an induction on the trace length.

\subsection{The proof-obligation record}

For each theorem package, the following record should be completed before
writing a large formal proof.

\begin{longtable}{@{}p{0.25\textwidth}p{0.67\textwidth}@{}}
\toprule
Field & Required content \\
\midrule
\endhead
Object presentation & Natural-number code or fixed evaluator schema.\\
Domain guard & Boolean validity, denominator positivity, boundedness,
nondegeneracy, or regularity assumptions.\\
Finite observations & Rational tables, residuals, grids, coefficient lists,
finite laws, traces, or local charts.\\
Precision transformer & Exact primitive-recursive function mapping output
precision to every requested input precision and truncation length.\\
Trace system & Initial state, step, primitive-recursive bound, and final
acceptance predicate.\\
QF invariant & Equation to one for a Boolean checker, including folded bounded
conditions.\\
Branch margins & Overlap or positive gap used for every decision affecting
partial operations.\\
Search/compactness certificate & Finite net, height bound, covering list,
isolating ball, or explicit finite candidate set.\\
Soundness bridge & Arithmetic implication from accepted finite data to the
stated observation.\\
Object-language target & Exact \code{Fm PTm} formula and intended
\code{Prf PRAAx} endpoint.\\
Native transport & Application of \code{pra_iff_native} only after the prior
field is complete.\\
\bottomrule
\end{longtable}

The archive includes this record as a JSON schema and a CSV difficulty matrix.
Its purpose is not bureaucratic: it prevents a theorem from being called
``formalized'' while its precision map, domain guard, or object-language target
is still implicit.

\subsection{Four theorem shapes}

The framework distinguishes four normal forms.

\begin{definition}[Finite exact theorem]
A finite exact theorem is
\[
 \forall p\;[\mathsf{Dom}(p)=1\to\mathsf{Check}(p,A(p))=1],
\]
where \(A\) and \(\mathsf{Check}\) are primitive recursive.
\end{definition}

\begin{definition}[Effective approximation theorem]
An effective approximation theorem is
\[
 \forall p\,\forall k\;
 [\mathsf{Dom}(p)\wedge\mathsf{InputContract}(p,\mu(p,k))
   \to \mathsf{OutputContract}(p,k,A(p,k))].
\]
All quantities in the matrix are finite rational statements.
\end{definition}

\begin{definition}[Certificate-relative existence theorem]
A certificate-relative theorem is
\[
 \forall p,c\;[\mathsf{Certificate}(p,c)=1
   \to \mathsf{Conclusion}(p,\mathsf{decode}(c))].
\]
The verifier is uniform.  A separate theorem may construct \(c\).
\end{definition}

\begin{definition}[Obstruction theorem]
An obstruction theorem shows that a stronger uniform selector or modulus
would realize a principle not validated by the intended constructive
semantics, or would compute a known noncomputable function.  The positive
replacement is stated in one of the preceding forms.
\end{definition}

These forms make status compositional.  A deep theorem can be incomplete only
at its certificate producer while its finite verifier, error propagation, and
formal bridge are finished and reusable.

\subsection{Composition}

Suppose \(F:X\to Y\) requests input precision \(\mu_F(p,j)\), and
\(G:Y\to Z\) requests \(Y\)-precision \(\mu_G(q,k)\).  The composite requests
\[
 \mu_{G\circ F}(p,k)=
 \mu_F\bigl(p,\mu_G(F_0(p),k)\bigr),
\]
where \(F_0(p)\) denotes the finite bound and domain metadata needed to choose
\(G\)'s precision.  If those bounds are themselves computed during the run,
concatenate the traces: first compute and certify the metadata, then run the
second module.  Pair the states, add the bounds, and conjoin the invariants by
a Boolean product.  Closure of primitive-recursive functions under
composition and primitive recursion gives a new certified transducer.

This observation turns the framework into a small category of proof-producing
finite transducers.  More importantly, it gives a practical library boundary:
rational arithmetic, finite folds, regular-name operations, robust decisions,
finite nets, and certificate checkers can be proved once and composed in many
packets.

\subsection{Proof-producing versus proof-checking modes}

There are two legitimate deployment modes.

\paragraph{Compiled mode.}
For a fixed theorem, the algorithm and all helper functions are represented by
fixed \code{PR} symbols.  Its trace is not an input; primitive recursion
computes it.  This mode gives the strongest uniform theorem and usually the
smallest final statement.

\paragraph{Certificate mode.}
A finite trace, isolating region, factorization, covering list, or analytic
constant table is input.  A small uniform checker validates it.  This mode is
appropriate when production is expensive, theorem-specific, or presently
unknown, and it is essential when no total universal producer can exist.

A mature development should often support both: certificate mode first, then a
producer theorem for the cases where an explicit construction is available.

\section{The finite-trace embedding theorem}
\label{sec:master}

\subsection{Trace systems}

All run-time parameters---the original input, requested precision, finite
bounds, and externally supplied certificates---may be paired into one natural
number \(u\).  A trace system consists of fixed primitive-recursive functions
\[
\begin{array}{ll}
 \mathsf{init}(u), & \mathsf{bound}(u),\\
 \mathsf{step}(u,i,s), & \mathsf{inv}(u,i,s)\in\{0,1\},\\
 \mathsf{good}(u)\in\{0,1\}, & \mathsf{accept}(u,s)\in\{0,1\}.
\end{array}
\]
Define the run by primitive recursion:
\[
 \mathsf{run}(u,0)=\mathsf{init}(u),\qquad
 \mathsf{run}(u,i+1)=\mathsf{step}(u,i,\mathsf{run}(u,i)).
\]
The intended proof obligations are the following quantifier-free formulas:
\begin{align}
 \mathsf{good}(u)=1&\to
   \mathsf{inv}(u,0,\mathsf{init}(u))=1, \tag{T0}\label{eq:T0}\\
 \mathsf{good}(u)=1\wedge \mathsf{inv}(u,i,s)=1&\to
   \mathsf{inv}(u,i+1,\mathsf{step}(u,i,s))=1, \tag{T1}\label{eq:T1}\\
 \mathsf{good}(u)=1\wedge
 \mathsf{inv}(u,\mathsf{bound}(u),s)=1&\to
   \mathsf{accept}(u,s)=1. \tag{T2}\label{eq:T2}
\end{align}

\begin{theorem}[Finite-trace reflection]
\label{thm:trace-reflection}
If \(\ipra\) proves \eqref{eq:T0}--\eqref{eq:T2}, then it proves
\[
 \forall u\;\bigl(\mathsf{good}(u)=1\to
 \mathsf{accept}(u,\mathsf{run}(u,\mathsf{bound}(u)))=1\bigr).
\]
\end{theorem}

\begin{proof}
Keep \(u\) free.  Apply quantifier-free induction to
\[
 A(i)\;:\equiv\;
 \mathsf{good}(u)=1\to
 \mathsf{inv}(u,i,\mathsf{run}(u,i))=1.
\]
This is quantifier free: all functions are fixed primitive-recursive symbols,
and implication is a propositional connective rather than a quantifier.  The
base case is \eqref{eq:T0} after rewriting \(\mathsf{run}(u,0)\).  For the
successor case, assume \(A(i)\) and \(\mathsf{good}(u)=1\).  The induction
hypothesis gives the invariant at \(\mathsf{run}(u,i)\); \eqref{eq:T1} gives
the invariant after one step; the defining equation for \(\mathsf{run}\)
identifies that state with \(\mathsf{run}(u,i+1)\).

Specialize the resulting universal statement to
\(i=\mathsf{bound}(u)\).  Apply \eqref{eq:T2} with
\(s=\mathsf{run}(u,\mathsf{bound}(u))\).  Finally introduce \(u\)
universally.  No induction formula contains an unbounded quantifier.
\end{proof}

\begin{remark}
The theorem is intentionally modest.  It does not say that every true
primitive-recursive assertion is provable.  It says that a supplied base case,
local invariant proof, and final implication compose into the expected bounded
algorithm theorem.  The mathematical work lies in choosing an invariant for
which \eqref{eq:T1} is actually derivable.
\end{remark}

\subsection{Finite query tables}

A represented input may have a global contract such as
\[
 \mathsf{Reg}_a(p)\;:\equiv\;
 \forall m,n\;R(a(p,m),a(p,n),m,n).
\]
Putting \(\mathsf{Reg}_a(p)\) inside the induction predicate would violate the
intended proof discipline.  Instead, the trace declares a finite query list
\(q(u,0),\ldots,q(u,L(u)-1)\).  Define \(\mathsf{GoodTable}(p,u)=1\) as the
Boolean conjunction of exactly the regularity, domain, and magnitude facts
needed about those samples.

\begin{lemma}[Finite instantiation]
\label{lem:finite-instantiation}
For every fixed query generator, \(\ipra\) proves
\[
 \mathsf{Reg}_a(p)\to\mathsf{GoodTable}(p,u)=1,
\]
provided each field of \(\mathsf{GoodTable}\) is an instance or a finite
algebraic consequence of the global contract.
\end{lemma}

\begin{proof}
The query list is finite and primitive recursive.  For each entry, eliminate
the universal quantifiers in \(\mathsf{Reg}_a(p)\) at the displayed numerical
indices.  The remaining rational equalities and inequalities are combined by
the Boolean-fold correctness lemma.  If later queries depend on earlier
answers, perform the same argument along the actual trace: at stage \(i\), the
query is a primitive-recursive term in the previous state, and universal
elimination permits that term.  The accumulator invariant says that all facts
queried so far have been inserted into the table.  This invariant is a Boolean
equation and is therefore quantifier free.
\end{proof}

This lemma explains how universally quantified name hypotheses can remain in
the surrounding context while the induction formula itself stays finite.

\subsection{Fix precision, then quantify}

\begin{theorem}[Fixed-precision compilation]
\label{thm:fixed-precision}
Let \(k\) be a requested output precision and let \(u=\angles{p,k,c}\) encode
all numerical inputs and finite certificates.  Suppose:
\begin{enumerate}
\item global hypotheses \(H(p)\) imply \(\mathsf{good}(u)=1\) by finite
      instantiation;
\item the trace system satisfies \eqref{eq:T0}--\eqref{eq:T2};
\item acceptance implies a finite output observation
      \(\mathsf{Out}(p,k,A(p,k))\).
\end{enumerate}
Then \(\ipra\) proves
\[
 H(p)\to\forall k\;\mathsf{Out}(p,k,A(p,k)).
\]
\end{theorem}

\begin{proof}
Keep both \(p\) and \(k\) free.  From \(H(p)\), derive the finite good table
for this particular \(k\).  Apply \cref{thm:trace-reflection}, then the final
soundness implication.  At no point is induction applied to a formula
containing \(\forall k\).  Once the finite derivation has been obtained with
\(k\) free, apply universal introduction to \(k\).
\end{proof}

This simple reordering resolves a large class of reported ``iPRA induction is
too weak'' concerns.  Frequently the problem is not the estimate but the
attempt to induct on its completed all-precision interpretation.

\subsection{Certificate reflection}

A computation need not be rerun inside the checker.  Let a certificate
\(c\) contain a state list \(s_0,\ldots,s_B\).  Define
\(\mathsf{TraceOK}(u,c)=1\) to assert
\[
 s_0=\mathsf{init}(u),\qquad
 s_{i+1}=\mathsf{step}(u,i,s_i)\ (i<B),\qquad
 \mathsf{inv}(u,i,s_i)=1\ (i\le B),
\]
and the final acceptance equation.  Every clause is checked by a finite fold.

\begin{corollary}[Uniform finite-certificate checker]
\label{cor:certificate-checker}
For a fixed trace language, \(\ipra\) proves
\[
 \mathsf{TraceOK}(u,c)=1\to\mathsf{Conclusion}(u,\mathsf{output}(c)).
\]
\end{corollary}

\begin{proof}
Use fold elimination to obtain every local transition and invariant clause.
The final clause directly supplies acceptance, or alternatively reconstruct
the run equality by induction on the list index and invoke
\cref{thm:trace-reflection}.  Both routes use only finite list induction with a
Boolean invariant.
\end{proof}

Certificate reflection is the right generic mechanism for externally produced
polynomial normalizations, isolating regions, factorization traces, finite
covers, and large numerical constant tables.

\subsection{Composition theorem}

\begin{theorem}[Composition of certified transducers]
\label{thm:composition}
The sequential composition of two \(\CFP\) transducers is a \(\CFP\)
transducer.  Its precision transformer, trace bound, and checker are primitive
recursive, and its soundness theorem is provable in \(\ipra\) from the two
component soundness theorems.
\end{theorem}

\begin{proof}
Encode a composite state as a phase bit together with the current component
state and the finite metadata passed between components.  In phase zero, run
the first transducer for its primitive-recursive bound.  Its final state
contains the output code and every certified magnitude/domain bound needed by
the second precision transformer.  Switch the phase bit, initialize the
second state, and run it for its bound.  The total bound is the sum of the two
bounds plus one switch step.

The composite invariant is a Boolean case split on the phase bit.  In phase
zero it is the first invariant.  At the switch it includes the first
acceptance predicate and the second input-good predicate.  In phase one it is
the second invariant together with the frozen metadata.  Each case is a
quantifier-free equation.  Apply finite-trace reflection.  The final soundness
implication is exactly the second component theorem after substituting the
certified first output.  Primitive-recursive closure under pairing,
composition, case distinction, addition, and primitive recursion supplies all
new symbols.
\end{proof}

The same construction handles finite acyclic proof networks and bounded
adaptive calls.  A network is flattened into one phase-indexed trace.

\subsection{Export to the native calculus}

Once the endpoint
\(\Gamma\vdash[\texttt{PRAAx}]\varphi\) has been constructed, the existing
repository theorem gives
\[
 \mathsf{NPrf}(\Gamma\text{ translated},\varphi\text{ translated}).
\]
No additional mathematical proof is required at this stage.  Conversely, the
native public/guarded calculi can serve as compact certificate targets when
their formula restrictions are convenient.  What must be avoided is invoking
bi-interpretability before there is any object-language proof to transport.

\section{Represented numbers, functions, and analytic operations}
\label{sec:represented}

\subsection{Integers and rationals}

An integer may be coded as a sign bit and magnitude, or by the standard zigzag
map.  A rational is a pair \((z,d)\) with \(z\in\Z\) and positive
\(d\in\N\).  Reduction to lowest terms is optional.  Equality and order are
checked by cross multiplication:
\[
 (z,d)=(w,e)\iff ze=wd,
 \qquad
 (z,d)\le(w,e)\iff ze\le wd,
\]
under \(d,e>0\).  Addition and multiplication use the unreduced formulas.
Absolute value, maximum, powers, finite sums, and finite products are
primitive recursive.

Clearing denominators is more than an implementation detail.  It converts the
correctness predicates for rational algebra into integer equations and
inequalities, which are directly suited to quantifier-free induction.  Every
rational theorem should therefore identify the common positive denominator or
state its cross-multiplied form.

\subsection{Regular names}

Let \(a(p,k)\in\Q\) be a fixed primitive-recursive evaluator.  Write
\(\delta_k=2^{-k}\).

\begin{definition}[Regular real name]
The parameter \(p\) is regular for \(a\) when
\[
 \mathsf{Reg}_a(p):\quad
 \forall k,l\quad
 \abs{a(p,k)-a(p,l)}\le \delta_k+\delta_l.
 \tag{R1}\label{eq:regular-name}
\]
\end{definition}

\begin{definition}[Equality and apartness]
Two fixed name families \(a,b\) are equal at parameters \(p,q\) when
\[
 \forall k,l\quad
 \abs{a(p,k)-b(q,l)}\le\delta_k+\delta_l.       \tag{R2}\label{eq:name-eq}
\]
They are apart with witness \(d\in\Q_{>0}\) when the finite-precision lower
bounds certify \(\abs{a-b}\ge d\); one convenient form is
\[
 \forall k,l\quad
 d\le \abs{a(p,k)-b(q,l)}+\delta_k+\delta_l.   \tag{R3}\label{eq:apart}
\]
Stronger one-sample variants are often easier to use.
\end{definition}

These are first-order formulas because \(a\) and \(b\) are fixed symbols and
all quantifiers range over natural precision indices.  They are schemas in the
evaluators, not quantification over a function space.

\begin{lemma}[Equality is transitive]
\label{lem:name-eq-trans}
Under regularity, \eqref{eq:name-eq} is an equivalence relation.
\end{lemma}

\begin{proof}
Reflexivity is \eqref{eq:regular-name}; symmetry is immediate.  For
transitivity, suppose \(a=b\) and \(b=c\).  For arbitrary \(k,l,j\),
\[
 \abs{a(k)-c(l)}
 \le \delta_k+\delta_l+2\delta_j.
\]
We must remove the extra \(2\delta_j\) without appealing to a real limit.
Assume the desired rational inequality fails.  Then the rational difference
\[
 r=\abs{a(k)-c(l)}-(\delta_k+\delta_l)
\]
is positive.  From the numerator and denominator of \(r\), bounded bit-length
search produces \(j\) with \(2\delta_j<r\).  The displayed upper bound then
contradicts the definition of \(r\).  This is finite rational arithmetic.  The
argument formalizes as a decidable case split on the cross-multiplied target
inequality and a primitive-recursive dyadic-below-positive-rational function.
\end{proof}

The same ``positive rational gap'' argument replaces every informal phrase
``let the error tend to zero'' used to prove an exact rational inequality
between names.

\subsection{Operations and precision transformers}

For every operation, the output algorithm and its error budget are recorded
together.  The following elementary contracts form a useful base library.

\begin{lemma}[Addition and finite sums]
If \(a_i,b_i\) approximate \(x_i,y_i\) within \(\eps_i\), then
\[
 \abs{(a_1+a_2)-(x_1+x_2)}\le\eps_1+\eps_2,
\]
and for a finite sum the errors add.  Thus a sum of \(m\) terms requested at
precision \(k\) may request each term at precision
\(k+\ceil{\log_2\max(1,m)}+1\).
\end{lemma}

\begin{proof}
Use the triangle inequality.  For the finite sum, induct on the accumulator
index with the invariant that its error is at most the sum of the processed
error radii.  The invariant is an integer inequality after denominators are
cleared.  The displayed precision shift makes the total less than
\(2^{-k}\).
\end{proof}

\begin{lemma}[Products]
\label{lem:product-error}
If \(\abs{x_i},\abs{y_i}\le B\) and
\(\abs{x_i-y_i}\le\eps\) for \(i<m\), then
\[
 \abs{\prod_{i<m}x_i-\prod_{i<m}y_i}
 \le m B^{m-1}\eps.
\]
\end{lemma}

\begin{proof}
Insert one changed factor at a time:
\[
 \prod x_i-\prod y_i=
 \sum_{j<m}\left(\prod_{i<j}x_i\right)(x_j-y_j)
                 \left(\prod_{j<i<m}y_i\right).
\]
Each summand is bounded by \(B^{m-1}\eps\).  The identity and bound are
proved by finite induction; no infinite product theorem is involved.
\end{proof}

\begin{lemma}[Reciprocal under apartness]
\label{lem:reciprocal}
Suppose \(\abs{x}\ge d>0\) and \(\abs{x-a}\le\eps\le d/2\).  Then
\(\abs{a}\ge d/2\) and
\[
 \abs{x^{-1}-a^{-1}}\le \frac{2\eps}{d^2}.
\]
\end{lemma}

\begin{proof}
The first claim is the reverse triangle inequality.  For the second,
\[
 \abs{x^{-1}-a^{-1}}=\frac{\abs{a-x}}{\abs{x}\abs{a}}
 \le \frac{\eps}{d(d/2)}.
\]
In the coded proof, division is absent from the invariant: multiply through by
the known positive denominators and by the square of the rational lower bound \(d\).
\end{proof}

\begin{lemma}[Polynomial Lipschitz bound]
\label{lem:poly-lip}
For \(P(X)=\sum_\alpha c_\alpha X^\alpha\), if every coordinate of \(x,y\)
has absolute value at most \(B\ge1\) and
\(\max_i\abs{x_i-y_i}\le\eta\), then
\[
 \abs{P(x)-P(y)}\le
 \left(\sum_\alpha\abs{c_\alpha}\,\abs{\alpha}\,B^{\abs{\alpha}-1}\right)\eta.
\]
\end{lemma}

\begin{proof}
For a monomial, telescope the product by replacing one factor at a time.
There are \(\abs{\alpha}\) summands, each with one difference factor and at
most \(\abs{\alpha}-1\) bounded factors.  Multiply by
\(\abs{c_\alpha}\) and sum over the finite coefficient list.  The proof is a
nested finite induction over monomials and factors.
\end{proof}

These lemmas give an automatic precision compiler for expression trees made of
finite sums, products, guarded inverses, and fixed polynomials.  Each node
returns both a value approximant and a certified magnitude bound for its
parent.

\subsection{Functions and continuity data}

A represented function is not merely an evaluator \(f(p,x,k)\).  It carries a
domain code and a modulus.  On a rational box, a modulus is a
primitive-recursive function \(\mu(k)\) satisfying
\[
 \dist(x,y)\le2^{-\mu(k)}\to
 \abs{f(x)-f(y)}\le2^{-k}.
\]
The arithmetic theorem uses rational approximants and slack, but the logical
shape is the same.  Composition uses \(\mu_f\circ\mu_g\).  A Lipschitz
constant supplies \(\mu(k)=k+\ceil{\log_2 L}+1\).

Derivatives should be represented by a local remainder contract, not by an
unqualified limit:
\[
 \abs{f(x+h)-f(x)-h g(x)}\le \abs{h}\,\omega(\abs{h}),
\]
with an explicit primitive-recursive modulus \(\omega\).  This is exactly the
form needed for a finite telescoping proof of the fundamental theorem on a
mesh.

\subsection{Series and integrals}

A series is represented by a finite partial sum and a tail bound
\(T(N)\).  To name its value at precision \(k\), compute a
primitive-recursive \(N(k)\) with \(T(N(k))\le2^{-k-c}\), evaluate the finite
sum, and allocate the remaining error to term approximation.  Regularity of
the resulting name follows by comparing two partial sums through the larger
cutoff and using the two tail bounds.

For an integral on \([a,b]\), use a rational mesh sum.  A Lipschitz constant
\(L\) gives
\[
 \abs{I_M-I_{M'}}\le L(b-a)^2(M^{-1}+{M'}^{-1}).
\]
Comparison on a common refinement gives this estimate.  The common refinement, double sum,
and error estimate are finite.  A primitive-recursive choice of \(M(k)\)
gives a regular name.  Changes of variable, integration by parts, and contour
integrals are admitted only with explicit local remainder or distance-to-pole
certificates; their proofs are finite telescoping calculations.

\subsection{Why no completeness axiom is needed}

The represented object is the evaluator and its regularity proof.  Operations
construct a new evaluator and prove its regularity.  The theory never asserts
that every abstract Cauchy sequence has a limit in a pre-existing real-number
sort.  Exact equality is the all-precision relation \eqref{eq:name-eq}; every
use of it is reduced to arbitrary finite precision and then introduced
universally.  This is sufficient for the analytic identities in the packets
that already carry effective approximation data.

\section{Robust control flow, gaps, and charts}
\label{sec:control}

\subsection{Why exact trichotomy is the wrong interface}

A finite algorithm can compare rationals exactly.  It cannot in general decide
whether two arbitrary regular names are equal, nor choose the sign of a name
that may be zero.  Many classical proofs use such a choice only to select a
formula, pivot, or coordinate chart.  The constructive repair is to make the
charts overlap, or to require a positive gap when a genuinely exact decision
is part of the conclusion.

\subsection{The overlapping branch lemma}

\begin{theorem}[Overlapping rational branch]
\label{thm:overlap-branch}
Let \(x\) be represented by a rational approximation \(a\) with
\(\abs{x-a}\le e\).  Let \(q\in\Q\), \(d>0\), and suppose \(4e\le d\).
A decidable comparison of \(a\) with \(q+d/2\) produces one of the valid
conclusions
\[
 x<q+d\qquad\text{or}\qquad x>q.
\]
The conclusions overlap on \((q,q+d)\).
\end{theorem}

\begin{proof}
If \(a\le q+d/2\), then
\[
 x\le a+e\le q+d/2+d/4<q+d.
\]
If \(a>q+d/2\), then
\[
 x\ge a-e>q+d/2-d/4>q.
\]
Only a rational comparison is performed.  In a pure name formulation, replace
strict inequalities by their rational-margin certificates.  All displayed
inequalities are checked after clearing denominators.
\end{proof}

This lemma is enough for overlapping analytic charts, small/large parameter
cases, approximate sign tests, and robust pivot selection.  It avoids the
false demand that the boundary case be assigned uniquely.

\subsection{Exact decisions from apartness}

\begin{theorem}[Exact branch under a supplied gap]
\label{thm:exact-gap}
Suppose a contract supplies \(d>0\) with
\(\abs{x-q}\ge d\).  If \(a\) approximates \(x\) within \(e<d/2\), then the
rational sign of \(a-q\) is nonzero and equals the sign of \(x-q\).
\end{theorem}

\begin{proof}
If \(x\ge q+d\), then \(a>q+d-e>q\).  If
\(x\le q-d\), then \(a<q-d+e<q\).  Conversely, an approximation equal to
\(q\) would be within \(e<d\) of \(x\), contradicting the gap.  The two cases
in the gap certificate may be encoded as a signed lower-bound witness or as a
Boolean comparison already supplied by the certificate.
\end{proof}

Apartness is therefore the correct input to inversion, normalization, root
ordering, exact sign, and uniqueness theorems.  Negative information
\(x\ne q\) is not interchangeable with a positive rational gap.

\subsection{Stable finite pivoting}

Let \(v_0,\ldots,v_{m-1}\) be rationally approximated candidate vectors, and
let \(r_i\) be a nonnegative residual or squared norm.  A finite maximum scan
returns the least index \(j\) attaining the largest sampled value.

\begin{lemma}[Robust pivot]
\label{lem:robust-pivot}
Suppose some exact residual satisfies \(r_i\ge d\), every sampled residual
\(\tilde r_i\) has error at most \(e\), and \(4e\le d\).  Then the selected
maximum satisfies \(r_j\ge d/2\).  Consequently any operation requiring a
positive lower bound may use the certified bound \(d/2\).
\end{lemma}

\begin{proof}
For the good index \(i\), \(\tilde r_i\ge d-e\).  Maximality gives
\(\tilde r_j\ge d-e\).  Hence
\(r_j\ge\tilde r_j-e\ge d-2e\ge d/2\).  The maximum scan invariant says that
the stored index dominates all processed samples and is a quantifier-free
Boolean fold.  No exact comparison of represented residuals is used.
\end{proof}

This is the generic core of stable column selection in matrix arguments,
choice of a nonzero polynomial coefficient, denominator selection in projective
coordinates, and finite chart selection on compact domains.

\subsection{Overlapping atlases}

A constructive atlas should be designed so that every point has at least one
chart with a quantitative denominator bound, and the chart-selection tests
have overlap.  For example, two formulas valid respectively when
\(c>-3/4\) and \(c<3/4\) cover all \(c\in[-1,1]\).  Approximate \(c\) to
error at most \(1/16\) and compare it to zero.  Either selected chart has a
fixed denominator margin.  On the overlap both are valid, so no proof of
unique chart selection is needed.

The same design applies to:
\begin{itemize}
\item quaternion charts for rotations;
\item affine charts of projective space;
\item inverse trigonometric formulas;
\item root formulas with competing dominant coefficients;
\item contour decompositions and local power-series disks.
\end{itemize}

Transition identities on overlaps should be finite polynomial/rational
certificates.  This replaces topological choice by a bounded rational decision.

\subsection{Projective and homogeneous output}

Normalization often creates an artificial discontinuity.  A nonzero kernel
vector has a continuous projective class even when no canonical unit vector is
available.  A polynomial root can be returned as an isolating disk or cluster
before an ordering is chosen.  A rotation can be returned as a unit quaternion
pair with vector part allowed to vanish at the identity.  A line or conic can
be represented by homogeneous coefficients with a positive maximum-norm
certificate, postponing division until a chart has been selected.

\begin{principle}[Delay normalization]
Whenever a theorem is invariant under scaling, sign, permutation, or another
finite symmetry, state and prove the invariant conclusion first.  Add a
canonical representative only under a separation or orientation certificate.
\end{principle}

This principle converts many apparent logical obstacles into ordinary finite
algebra.

\subsection{Partial operations as certified transducers}

A guarded inverse transducer takes a name \(x\), a positive rational \(d\),
and a proof that \(\abs{x}\ge d\).  Its precision transformer requests enough
input accuracy that the rational approximation is at least \(d/2\) away from
zero; it then uses rational division.  The trace records the comparison and
cross-multiplied error inequality.  Square root, logarithm, matrix inversion,
and normalization have the same shape:

\[
\text{domain/apartness bit}
\;\longrightarrow\;
\text{precision request}
\;\longrightarrow\;
\text{finite rational operation}
\;\longrightarrow\;
\text{error certificate}.
\]

The operation itself remains total on all codes, but only the guarded theorem
assigns it the intended mathematical meaning.

\section{Finite nets, search, and exact objects from nested certificates}
\label{sec:existence}

\subsection{Effective finite nets}

\begin{definition}[Certified net presentation]
A bounded metric-object family has a certified net presentation when there are
primitive-recursive functions \(N(k)\) and \(\mathsf{net}(k,i)\),
\(i<N(k)\), together with a theorem that every valid input object lies within
\(2^{-k}\) of some listed point.  The coverage theorem may compute the index
from a finite approximation, or accept it as a certificate.
\end{definition}

For rational boxes, simplices, circles with overlapping charts, finite
probability simplices, and coefficient-bounded polynomials, the standard grids
are primitive recursive.  Abstract compactness is not assumed; the finite net
and its coverage modulus are data.

\begin{theorem}[Finite-net transfer]
\label{thm:net-transfer}
Let \(f\) have modulus \(\mu\) on a certified net domain.  If \(x\) is a
valid named point and \(i<N(\mu(k))\) satisfies
\(\dist(x,\mathsf{net}(\mu(k),i))\le2^{-\mu(k)}\), then
\[
 \abs{f(x)-f(\mathsf{net}(\mu(k),i))}\le2^{-k}.
\]
The corresponding arithmetic statement is provable in \(\ipra\).
\end{theorem}

\begin{proof}
This is the defining implication of the modulus instantiated at the finite net
point.  In the arithmetic proof, approximate both sides to a smaller error,
use the finite coverage certificate, and clear denominators.  All queried
precisions are computed from \(k\), so fixed-precision compilation applies.
\end{proof}

\subsection{Approximate existence by bounded minimization}

Let a nonnegative residual \(R(y)\) vanish exactly on desired solutions and
have a known continuity modulus.  Scan a finite net and return the least index
with minimal sampled residual.

\begin{theorem}[Approximate witness from a net]
\label{thm:approx-witness}
Suppose a valid named \(x\) satisfies \(R(x)=0\), and
\(\abs{R(y)-R(z)}\le L\dist(y,z)\) on the domain.  On a net of radius \(r\),
the bounded minimizer \(y_*\) satisfies
\[
 R(y_*)\le Lr+2e,
\]
where \(e\) is the error in each sampled residual.  The algorithm and bound
are primitive recursive.
\end{theorem}

\begin{proof}
Choose a net point \(y_i\) within \(r\) of \(x\).  Then
\(R(y_i)\le Lr\).  Its sampled residual is at most \(Lr+e\).  The selected
sample is no larger, and converting its sample back to the exact residual adds
another \(e\).  The scan invariant records the least processed index with
minimal rational sample.  It is quantifier free after a bounded comparison
fold.
\end{proof}

This theorem is often the strongest uniform constructive replacement for an
exact minimizer, eigenvector, axis, or zero.  It gives a witness at every
requested accuracy without asserting that the selected witnesses converge.

\subsection{Coherent nested certificates}

An exact represented object can be obtained when the finite witnesses are
coherent.

\begin{definition}[Nested isolating certificate]
A nested certificate is a primitive-recursive sequence of rational balls
\(B_k=B(c_k,r_k)\) such that
\[
 r_k\le2^{-k},\qquad B_{k+1}\subseteq B_k,
\]
and a finite checker proves that every \(B_k\) has the desired existence
property.  If uniqueness is needed, a separation certificate says that at
most one solution lies in \(B_k\).
\end{definition}

\begin{theorem}[Nested balls produce a regular name]
\label{thm:nested-name}
The center sequence of a nested certificate is a regular name after a fixed
precision shift.  Every property stable under the certified residual bounds
holds of that name.  If each ball contains at most one solution, the represented
solution is unique in the name equality relation.
\end{theorem}

\begin{proof}
For \(l\ge k\), nesting puts \(c_l\) in \(B_k\), hence
\(\dist(c_k,c_l)\le r_k+r_l\le2^{-k}+2^{-l}\).  For arbitrary \(k,l\), use
\(m=\min(k,l)\) and the same estimate, or symmetrize.  This is exactly the
regular-name condition.  A residual bound tending with the displayed dyadic
rate yields the all-precision exact property by fixed-precision compilation.
If two represented solutions lie in every isolating ball, the uniqueness
clause at any requested precision bounds their distance; the positive-rational
gap argument proves name equality.
\end{proof}

Applications include polynomial roots via isolating disks, solutions of
well-conditioned implicit equations, normalized eigenvectors under a spectral
gap, and limits of contractive iterations.  The crucial input is not the word
``compact'' but the explicit refinement and isolation checker.

\subsection{Approximate minimizers without an exact minimizer theorem}

On a finite net, a minimum always exists by bounded scan.  With a uniform
continuity modulus, these minima form an approximate-infimum name:
\[
 m_k=\min_{i<N(\mu(k))} f(\mathsf{net}(\mu(k),i)).
\]
Comparing two net levels through a refinement or through the true pointwise
infimum bounds \(\abs{m_k-m_l}\).  Thus the infimum value can often be named
without choosing a minimizing point.  An exact minimizer requires additional
coherence, a compact choice principle, or a quantitative uniqueness condition.
This distinction is useful in optimization, spectral estimates, and extremal
geometry.

\subsection{Finite covers}

Suppose every net point has a certified neighborhood on which one member of a
finite family of local assertions holds.  A Lebesgue-number-style result is
constructive only when the neighborhood radii have a positive finite minimum
that is actually computed.  The checker scans the finite list, verifies every
radius, takes their minimum, and proves coverage by the net-transfer theorem.
There is no appeal to an abstract open-cover compactness theorem.

\subsection{Search bounds and finite lists}

Let \(P(a,x)\) be decidable and let \(h(x)\) be a primitive-recursive height
such that all codes of height at most \(B\) have a primitive-recursive finite
enumeration.

\begin{theorem}[Height bound and complete list]
\label{thm:height-list}
The following data are primitively and intuitionistically interconvertible:
\begin{enumerate}
\item a bound \(H(a)\) with \(P(a,x)\to h(x)\le H(a)\);
\item a finite list \(L(a)\) that is sound and complete for \(P(a,\cdot)\).
\end{enumerate}
\end{theorem}

\begin{proof}
Given \(H(a)\), enumerate all codes of height at most \(H(a)\) and retain
exactly those passing the decidable predicate.  Bounded-filter correctness is
a Boolean-fold induction.  Completeness follows from the height implication.

Given \(L(a)\), take the maximum of the heights of its entries, with zero for
the empty list.  If \(P(a,x)\), completeness gives membership of \(x\) in the
list; finite maximum correctness gives the height bound.  No excluded middle
or choice is used.
\end{proof}

A bound on \emph{the number} of solutions is absent from this equivalence.  It
does not locate the solutions and supplies no stopping criterion; a precise
computability obstruction is proved in \cref{thm:cardinality-not-height}.

\section{Reusable domain libraries}
\label{sec:domains}

The framework is most useful when it becomes a small collection of shared
object-language libraries.  This section describes the interfaces and the
proof reductions they support.

\subsection{Finite algebra and polynomial certificates}

A multivariate integer polynomial is a finite list of coefficient/exponent
pairs in a canonical order.  Addition and multiplication merge and convolve
those lists.  Normalization combines equal monomials and removes zero
coefficients.  Evaluation is a finite fold.

\begin{proposition}[Polynomial certificate reflection]
\label{prop:poly-reflection}
There is a primitive-recursive checker
\(\mathsf{PolyId}(P,Q,c)\) such that \(\ipra\) proves: if \(c\) is an accepted
normalization certificate for \(P-Q\), then
\[
 \forall \mathbf x\quad P(\mathbf x)=Q(\mathbf x).
\]
The same method checks rational identities after multiplying by a certified
positive common denominator.
\end{proposition}

\begin{proof}
The certificate lists every normalization step: sorting, combining equal
monomials, and deleting zeros.  The checker verifies local list equations.
Induction on the certificate length proves that evaluation is preserved.
Acceptance says the final list is empty, whose evaluation is zero.  The
induction invariant is equality of two integer evaluations and is
quantifier-free.  Variables \(\mathbf x\) remain free throughout and are
universally introduced at the end.
\end{proof}

This library handles ring identities, determinant expansions, adjugate
identities, trigonometric polynomial identities in \(\sin\) and \(\cos\), and
the finite coefficient parts of power-series proofs.  It also separates two
checks that are often conflated: a computer algebra system may produce the
certificate, while the small arithmetic checker proves its soundness.

\paragraph{Algebraic roots.}
A root should be represented first by an isolating rectangle or disk, a root
count, and a refinement procedure.  Polynomial arithmetic, interval bounds,
and a finite argument-principle or subdivision certificate validate each
region.  Repeated roots are represented by clusters with multiplicity data.
An ordered list of individual roots is added only under pairwise separation.
This resolves the branch discontinuity at collisions without weakening the
algebraic theorem.

\subsection{Linear algebra and geometry}

Matrices over coded rationals are finite lists.  Matrix multiplication,
transpose, determinant, minors, adjugate, dot products, and cross products are
finite polynomial operations.  Fixed-dimension identities are direct
polynomial certificates; variable-dimension identities use fold invariants.

\paragraph{Nonsingular systems.}
Cramer's rule and adjugate inversion are constructive under a determinant
apartness witness.  The transducer carries a lower bound on
\(\abs{\det A}\), requests sufficient entry precision, and proves the inverse
error by \cref{lem:reciprocal,lem:product-error}.  Gaussian elimination is
also possible, but every pivot choice must use a robust lower-bound scan.

\paragraph{Singular systems and eigenvectors.}
An exact nonzero kernel vector is not a continuous function of a singular
matrix near rank changes.  The unconditional output should be one of:
\begin{itemize}
\item a projective kernel certificate;
\item a residual-small unit vector at requested precision;
\item an exact normalized vector under a supplied minor or spectral gap;
\item a coherent isolating subspace projector rather than a basis.
\end{itemize}
Finite-net minimization of \(\norm{Av}^2\) on a rational sphere chart gives
the approximate version.  A gap promotes it to a stable exact name.

\paragraph{Euclidean and projective geometry.}
Points, lines, circles, conics, and incidences are homogeneous coordinate
codes.  Classical incidence theorems often reduce to determinant or polynomial
identities after nondegeneracy denominators are recorded.  Power-of-a-point,
Pappus, Ceva, Menelaus, and many circle theorems fit this pattern.  The library
should expose determinant certificates, signed-area predicates, and chart
transitions, not a large opaque geometry ontology.

\paragraph{Trigonometry and spherical geometry.}
Use regular names for \(\sin\) and \(\cos\) built from finite series or a
circle parameter, and state identities first as polynomial equations in a
pair \((c,s)\) satisfying \(c^2+s^2=1\).  Division-based tangent or half-angle
forms require denominator gaps.  Spherical theorems should prefer vector dot
and triple products; inverse angles are added via overlapping charts and
separation certificates.

\subsection{Real and complex analysis}

The shared analytic library should have four layers.

\begin{enumerate}
\item \textbf{Regular names and interval arithmetic.}  This is the numerical
      kernel of \cref{sec:represented}.
\item \textbf{Finite sums, products, and series tails.}  Every named special
      function exports a truncation algorithm, magnitude bound, and tail
      modulus.
\item \textbf{Mesh calculus.}  Integration, differentiation, and changes of
      variable use finite mesh sums plus local remainder moduli.
\item \textbf{Local analytic data.}  A holomorphic chart contains coefficient
      names and a majorant on a larger rational disk.  Cauchy estimates and
      termwise operations become finite tail calculations.
\end{enumerate}

A contour is a finite polygonal or piecewise rationally parametrized path.  A
residue theorem instance needs a positive distance from the path to every
pole, isolating disks for enclosed poles, local Laurent coefficient
certificates, and a finite winding-number computation.  These are exactly the
inputs that a classical diagram tends to suppress.

A uniform convergence theorem must include a primitive-recursive modulus.  A
pointwise statement without a modulus generally does not support interchange
of limit, integral, or derivative in \(\ipra\).  Dominated convergence is
replaced by finite truncation plus a tail or uniform-integrability modulus.

\subsection{Elementary and Diophantine number theory}

Bounded trial division, Euclid's algorithm, extended gcd, modular arithmetic,
finite factorization, valuations below a supplied bound, and Chinese remainder
construction are primitive-recursive finite algorithms.  Their proofs use
state invariants such as
\[
 r_i=A u_i+B v_i,
 \qquad
 \gcd(r_i,r_{i+1})=\gcd(A,B),
\]
encoded by integer equations and a bounded decreasing measure.

Diophantine geometry adds a search interface:
\begin{itemize}
\item a decidable point predicate;
\item a height and enumeration of bounded-height codes;
\item a height, descent, Runge, or covering certificate;
\item a finite verifier proving that the certificate applies to the input.
\end{itemize}
The generic descent theorem can then be fully embedded even when the production
of a global height certificate remains open.

\subsection{Analytic number theory}

Analytic-number-theory proofs should be divided into a finite arithmetic core
and an analytic input contract.

\begin{definition}[Analytic input certificate]
An analytic input certificate is a finite code containing explicit constants,
regions, exceptional cases, truncation lengths, and rational inequalities
sufficient for the later finite estimate.  Its checker validates every finite
algebraic and numerical condition.
\end{definition}

For a zero-free region, for example, the certificate records the region
parameters, a finite cover or local argument data, bounds on coefficients and
tails, and any exceptional zero with a separation estimate.  The downstream
prime-counting proof is a certified transducer from this data to an explicit
threshold.  The theorem is valuable even before a uniform certificate producer
is known: it states exactly which quantitative input is missing and prevents
an ineffective existence statement from entering unnoticed.

\subsection{Finite and effective probability}

A finite rational law is a finite list of rational values and natural weights
with positive total weight.  Expectation is an exact finite weighted sum.  A
product law is a mixed-radix enumeration; independence is an identity of
finite sums rather than an axiom.

Moments, characteristic functions at rational arguments, convolution, and
events with rational thresholds are finite calculations.  A Berry--Esseen or
Lindeberg proof is then a chain of polynomial/trigonometric approximation
certificates, finite product telescoping, and rational smoothing estimates.

An infinite or continuous law is admitted only through finite rational
approximants with explicit tail and quantization moduli.  For weak convergence,
a bounded-Lipschitz or characteristic-function metric is usually more robust
than CDF comparison.  Transfer to CDFs requires a continuity or
anti-concentration modulus near the threshold.  This cleanly separates the
finite probabilistic theorem from the presentation of the input law.

\subsection{Combinatorics and topology}

Finite combinatorics is a natural fit: graphs, words, partitions, tableaux,
simplicial complexes, and finite colorings are codes, while recurrences and
injections are bounded traces.  Generating-function proofs should be compiled
to coefficient identities before analytic notation is introduced.

Constructive topology uses subdivisions and parity certificates.  Sperner,
Brouwer-style approximate fixed points, polygonal Jordan separation, and
finite-dimensional degree arguments can be stated on rational grids with
explicit moduli.  Exact points arise from nested cells only when coherence is
proved.  Dimension and invariance results require careful negative statements,
but their local combinatorial certificates still fit finite-trace reflection.

\section{Resolution of the principal reported difficulties}
\label{sec:resolutions}

This section applies the common framework to the difficulties identified in
the repository.  ``Resolved'' means that the logical and representation
problem has a precise constructive replacement and a clear embedded-proof
route.  It does not mean that every deep theorem-specific estimate has been
proved.

\subsection{Unbounded precision and quantifier-free induction}

\paragraph{Problem.}
A report defines equality or convergence by a formula quantified over all
precisions, then attempts to use that formula as an induction invariant.  The
invariant is not quantifier free.

\paragraph{Resolution.}
Use \cref{thm:fixed-precision}.  Keep requested precision \(k\) free.  Compute
a finite query table containing only the samples needed at that \(k\).  Encode
all bounded conditions as Boolean folds.  Prove the finite accumulator theorem
by \code{PRA.qfInd}.  Introduce \(k\) universally at the end.  Global name
regularity is used only to derive the finite \(\mathsf{GoodTable}\) by
universal elimination.

\paragraph{Consequence.}
Most explicit series, mesh integral, finite-product, regular-name operation,
and finite probability proofs need no stronger induction.  A genuine problem
remains only when no primitive-recursive modulus or finite query bound is
available.

\subsection{Euler axes, eigenvectors, and discontinuous selectors}

\paragraph{Problem.}
At the identity rotation every unit vector is an axis.  A uniform exact unit
axis requires a discontinuous choice; the Euler--Rodrigues packet gives an
LLPO reduction.

\paragraph{Resolution.}
Use a four-tier interface:
\begin{enumerate}
\item unconditional exact homogeneous quaternion data, allowing zero vector
      part at the identity;
\item unconditional axis-and-angle approximation at every requested precision,
      with no claim that the axes converge;
\item exact projective or normalized axis under a positive separation of the
      rotation from the identity or under a nonzero-minor certificate;
\item exact rational axis for rational proper orthogonal matrices, where a
      finite nonzero cofactor or column can be found by exact comparison.
\end{enumerate}
Robust pivoting proves the separated case; polynomial certificates prove the
Rodrigues identities.  The unrestricted selector remains correctly marked
obstructed.

This pattern extends to eigenvectors at multiple eigenvalues, singular kernel
vectors, root ordering at collisions, and orientation choices.

\subsection{Dirichlet and Kronecker endpoints}

\paragraph{Problem.}
Closed sharp endpoints or density from negative irrationality ask the
arithmetic to decide a boundary or extract positive information from a double
negation.  The corresponding packet derives LLPO, LPO, or a Markov principle
from several unrestricted formulations.

\paragraph{Resolution.}
State strict approximation with a displayed rational margin, or accept a
positive irrationality/separation modulus.  Finite pigeonhole and grid
arguments then return an exact rational witness.  For matrix Kronecker, use a
finite compatibility certificate and an approximation-or-small-relation
alternative.  If the separation modulus rules out the relation branch, the
approximation branch follows.

The framework explanation is uniform: finite search can decide rational boxes;
it cannot decide whether a represented boundary value is exactly attained.
An overlap or gap removes the discontinuity.

\subsection{Faltings and Siegel: finiteness versus a stopping rule}

\paragraph{Problem.}
A classical theorem says that a rational- or integral-point set is finite, or
gives a bound on its cardinality.  A constructive complete list requires a
height cutoff.  The Faltings and Siegel packets correctly do not claim a
general cutoff.

\paragraph{Resolution.}
Build and embed the following modular theorem:
\[
 \mathsf{HeightCertificate}(C,H)=1
 \to
 \mathsf{CompleteList}(C,\mathsf{Search}(C,H)).
\]
The checker verifies the curve, certificate hypotheses, and all finite search
steps.  Conversely, a complete list supplies a maximum height by
\cref{thm:height-list}.  Thus the exact missing mathematical object is a
height-certificate producer, not finite search or list verification.

For Runge families, Fermat descents, or other explicit cases, the producer is
available and the whole theorem can be embedded.  For general Faltings or
Siegel, the verifier and equivalence theorem remain useful while the producer
is honestly open.  A cardinality estimate is recorded as different data and
is never used as a cutoff.

\subsection{Mordell--Weil and covering representatives}

\paragraph{Problem.}
A finite quotient or Selmer computation need not provide rational points
representing every relevant coset.  Descent needs actual representatives to
write \(P=2Q+R_i\).

\paragraph{Resolution.}
Define a covering certificate consisting of a finite rational point list
\(R\), validity proofs for its entries, and a checked covering relation or a
finite trace deriving the needed decomposition in the target family.  From
this certificate, a bounded rational-halving scan constructs a selector and
the generic height descent is finite-trace reflection.  The descent theorem is
therefore unconditional \emph{relative to a checked covering certificate}.
Producing such certificates uniformly for arbitrary curves remains separate.

This factorization improves the theorem interface: an uncheckable statement
``assume representatives exist'' is replaced by a finite object whose use is
fully verified.

\subsection{Ineffective constants in prime distribution}

\paragraph{Problem.}
An asymptotic theorem needs a threshold \(N(A,q)\) before the later variable
\(x\).  A classical proof may only assert existence of a constant or may hide
an exceptional-zero dependence.  Choosing a new threshold after seeing each
finite \(x\) proves the wrong quantifier order.

\paragraph{Resolution.}
Separate two modules:
\begin{align*}
 \mathsf{ProduceAnalyticCertificate}(D)&=c,\\
 \mathsf{VerifyAndPropagate}(D,c,A)&=N.
\end{align*}
The second is a finite arithmetic theorem and should be embedded first.  Its
input code contains every constant, exceptional object, separation margin, and
local estimate used downstream.  Its conclusion has the correct order
\[
 \forall A\;\exists N\;\forall x\ge N\;E(D,A,x).
\]
When a primitive-recursive producer is later proved for a class of data, the
composition theorem yields the unconditional result for that class.

This resolves the formal architecture and prevents accidental ineffectivity.
It does not manufacture the missing number-theoretic estimate.

\subsection{Ambient Lean versus embedded proof}

\paragraph{Problem.}
Many packets provide compilable or plausible Lean theorems about native data
but stop before an object-language derivation.

\paragraph{Resolution.}
Use a three-stage artifact policy:
\begin{enumerate}
\item semantic reference implementation and exact tests;
\item object-language specification with fixed \code{PR} symbols, formulas,
      and QF invariants;
\item actual \code{Prf PRAAx} construction, followed by
      \code{pra_iff_native} and an axiom audit.
\end{enumerate}
A status table must mark the highest completed stage for every main theorem.
The generic trace-reflection module reduces stage three to a small number of
shared proofs rather than reimplementing induction plumbing in each packet.

\subsection{Root choice and the fundamental theorem of algebra}

\paragraph{Problem.}
A classical FTA statement may return an unordered multiset of roots, while a
formal development asks for a continuously ordered root list.  Collisions make
such ordering discontinuous; analytic existence may also be written without a
computable radius or winding certificate.

\paragraph{Resolution.}
Return finite root clusters with rational isolating disks and multiplicities.
Refine each disk by a checked subdivision procedure.  The centers yield names
by \cref{thm:nested-name}.  Split and order clusters only under a positive
separation certificate.  Polynomial reconstruction from clusters is a finite
coefficient certificate.  Thus the algebraic content survives without a
uniform labeling choice at collisions.

\subsection{Central limits and arbitrary probability spaces}

\paragraph{Problem.}
The textbook theorem quantifies over measures and real random variables, and
CDF events may involve irrational thresholds or discontinuities.

\paragraph{Resolution.}
Use finite rational laws as the exact core.  State the finite triangular-array
estimate at rational raw thresholds.  Present more general laws by finite
rational approximants with primitive-recursive moment-tail and quantization
moduli.  Prove convergence first in a robust metric; transfer to CDFs with an
explicit continuity or anti-concentration bound.  The unrestricted qualitative
theorem becomes a schema for effective presentations, not a literal
higher-type sentence.

\section{Sharp boundaries and impossibility arguments}
\label{sec:barriers}

A general framework is trustworthy only if it explains what it cannot encode.
The following barriers show that several recurring difficulties are
mathematical, not deficiencies of notation.

\subsection{No total universal primitive-recursive evaluator}

\begin{theorem}[Diagonal obstruction]
\label{thm:no-universal-pr}
There is no binary primitive-recursive function \(U(e,x)\) such that every
unary primitive-recursive function \(f\) has an index \(e\) satisfying
\(f(x)=U(e,x)\) for all \(x\).
\end{theorem}

\begin{proof}
If such \(U\) existed, then
\[
 g(n)=U(n,n)+1
\]
would be primitive recursive.  By universality, choose \(e\) with
\(g(n)=U(e,n)\) for every \(n\).  At \(n=e\),
\[
 U(e,e)=g(e)=U(e,e)+1,
\]
a contradiction.
\end{proof}

Consequently, a first-order theorem cannot receive an arbitrary PR function
code and evaluate it by one total PR symbol.  Fixed evaluator schemas and
finite trace verifiers are not a temporary workaround; they are the maximal
uniform interface compatible with primitive recursiveness.

\subsection{A cardinality bound does not yield a height bound}

\begin{theorem}[No uniform conversion from ``at most one'' to a cutoff]
\label{thm:cardinality-not-height}
There is a primitive-recursive decidable predicate \(P(e,t)\) such that for
every \(e\) at most one \(t\) satisfies \(P(e,t)\), but there is no
primitive-recursive function \(H(e)\) satisfying
\[
 P(e,t)\to t\le H(e)
\]
for all \(e,t\).
\end{theorem}

\begin{proof}
Let \(P(e,t)\) say that Turing machine \(e\), on blank input, halts for the
first time exactly at step \(t\).  Bounded deterministic simulation is
primitive recursive, so \(P\) is primitive-recursive and decidable.  A machine
has at most one first halting time.

If a primitive-recursive \(H\) bounded every possible first halting time, the
halting problem would be decidable: simulate machine \(e\) for \(H(e)\) steps.
If it has not halted by then, the implication says it never halts.  This
contradicts undecidability of the halting problem.
\end{proof}

Thus even a uniform cardinality bound of one carries no search-stopping
information.  The obstacle in the Faltings and Siegel packets is exactly the
missing location/height data, not the finite filtering algorithm.

\subsection{Exact sign selection and LLPO}

One standard form of LLPO says: for every binary sequence with at most one
entry equal to one, either all even entries are zero or all odd entries are
zero.  From such a sequence \(\alpha\), construct the regular real name
corresponding to
\[
 x=\sum_{n\ge0}(\alpha(2n)-\alpha(2n+1))2^{-n-1}.
\]
The at-most-one premise gives an explicit tail modulus.  If a uniform
procedure always decides \(x\ge0\) or \(x\le0\), then:
\begin{itemize}
\item \(x\ge0\) rules out an odd-position one;
\item \(x\le0\) rules out an even-position one.
\end{itemize}
When the sequence is all zero, both conclusions are valid.  Hence real sign
selection realizes LLPO.

The exact-axis and sharp-endpoint obstruction arguments in the inspected
packets reduce their selectors to this sign principle.  Replacing the real
encoding or choosing a different normal form cannot remove the obstruction.
Overlapping branches avoid it because they do not decide which side of an
exact boundary contains zero.

\subsection{Negative inequality is not positive apartness}

Intuitionistically, \(\neg(x=0)\) need not provide a rational \(d>0\) with
\(\abs{x}\ge d\).  Algorithms for inversion, normalization, or density need
the latter.  A principle converting every negative nonzeroness proof into a
positive modulus is a Markov-style principle for represented numbers.  The
framework therefore treats apartness witnesses as data and never derives them
from bare negation unless the input class has a specific finite argument.

For rationals and algebraic numbers with an isolating certificate, positive
apartness can be computed.  For arbitrary regular names, it cannot.

\subsection{Ineffective existence cannot be repaired by a verifier}

A finite verifier may prove that any supplied constant table implies an
asymptotic estimate.  This does not imply a primitive-recursive function
producing such a table.  The difference mirrors proof checking versus theorem
search: local verification may be elementary even when uniform production is
open or noncomputable.

Quantifier order detects the issue.  The useful statement is
\[
 \forall D\,\forall A\;\exists N\;\forall x\ge N\;E(D,A,x),
\]
with \(N\) produced before \(x\).  A family of proofs of
\(\forall D,A,x\;\exists N_x\le x\;E(D,A,x)\) is not a substitute.  The
finite-trace framework preserves the former order by making the threshold an
output of a certificate producer.

\subsection{Semantic truth is not automatic iPRA provability}

A primitive-recursive function can have true arithmetic properties whose
proofs require induction or consistency strength beyond a particular weak
theory.  Therefore the slogan ``the algorithm is PR, so iPRA proves it'' is
false.  The framework always asks for the local invariant derivations
\eqref{eq:T0}--\eqref{eq:T2}.  The finite-trace theorem composes proofs; it
does not create missing mathematical invariants.

\subsection{Maximal positive statement principle}

When an unrestricted theorem meets one of these barriers, the goal is not to
stop at ``impossible.''  One should determine the maximal positive family:
\begin{itemize}
\item approximate witness at every precision;
\item exact invariant or projective data;
\item exact witness under a positive gap;
\item certificate-relative theorem;
\item theorem for rational or finitely presented inputs;
\item obstruction showing why the omitted uniform selector is stronger.
\end{itemize}
The Euler--Rodrigues and Dirichlet--Kronecker packets already exemplify this
pattern.  The \(\CFP\) framework makes it a standard deliverable for every
future packet.

\section{Lean formalization architecture}
\label{sec:lean}

\subsection{Existing interfaces}

The implementation should extend, not duplicate, the following inspected
modules.

\begin{longtable}{@{}p{0.37\textwidth}p{0.55\textwidth}@{}}
\toprule
Module/declaration & Role \\
\midrule
\endhead
\code{ConciseIPRA/PR.lean}: \code{PR}, \code{PTm} & Fixed primitive-recursive
function symbols and arithmetic terms.\\
\code{PRAAx} & Defining equations, successor axioms, and quantifier-free
induction.\\
\code{PRA.precZero}, \code{PRA.precSucc} & Rewriting equations for a concrete
primitive-recursive trace runner.\\
\code{PRA.qfInd} & The main invariant rule, with an explicit \code{QF} proof.\\
\code{ConciseIPRA/Native.lean}: \code{NFm}, \code{NPrf} & Native scalar
formulas and proofs.\\
\code{ConciseIPRA/NativeBiInterpretation.lean}: \code{pra_iff_native} &
Faithful export of a completed \code{PRAAx} derivation.\\
\code{ConciseIPRA/NativePublic.lean}: \code{PGPrf}, public gating & Optional
compact, substitution-stable public certificate language.\\
\bottomrule
\end{longtable}

The core arithmetic formalization should target \code{PRAAx}.  This keeps
rational coding and Boolean invariants readable.  Native translation is the
last step.

\subsection{The trace-system data}

After pairing all numerical parameters into one code \code{u}, a Lean
meta-level record can hold the fixed object-language symbols:

\begin{lstlisting}[style=leanstyle]
structure PRTraceSystem where
  good   : PR 1
  init   : PR 1
  bound  : PR 1
  step   : PR 3   -- arguments: i, previousState, u
  inv    : PR 3   -- arguments: i, state, u
  accept : PR 2   -- arguments: state, u

def PRTraceSystem.run (T : PRTraceSystem) : PR 2 :=
  .prec T.init T.step
\end{lstlisting}

The order matches the repository's primitive-recursion convention: the
recursion argument comes first, and the step receives the predecessor, prior
value, and parameters.  Pairing can reduce larger arities to this interface.

Define an object-language predicate \code{IsOne t} as the equation of a
Boolean term with the chosen code for one.  For a concrete \code{T}, construct
terms
\[
 \mathsf{Run}(i,u),\quad \mathsf{Good}(u),\quad
 \mathsf{Inv}(i,u),\quad \mathsf{Accept}(u).
\]
Do not create a second Boolean or rational encoding if the repository already
has one suitable for reuse.

\subsection{Embedded proof of trace reflection}

The dependency-ordered Lean proof is:

\begin{enumerate}
\item Prove simplification lemmas for \code{T.run} by instantiating
      \code{PRA.precZero} and \code{PRA.precSucc}.
\item Define the formula
\[
 \phi(u,i)=\mathsf{Good}(u)\to\mathsf{Inv}(i,u).
\]
      Prove \code{phi.QF} structurally.  Every bounded assertion inside
      \code{Inv} must already be a Boolean-fold equation.
\item Supply the base derivation from the concrete module's base obligation.
\item Under \code{phi :: Gamma.map Fm.shift}, use the induction hypothesis,
      the concrete step obligation, and the successor equation for
      \code{T.run} to derive \code{phi} at successor.
\item Invoke \code{PRA.qfInd phi hQF h0 hS}.
\item Eliminate the universal at \code{T.bound u}, apply final soundness, and
      introduce \code{u} universally.
\end{enumerate}

The generic theorem should take the base, step, and final derivations as fields
of a record.  This makes theorem-specific files provide only local arithmetic
lemmas.

\subsection{Boolean folds and bounded search}

The first supporting module should implement object-language symbols and
proofs for:
\begin{itemize}
\item pair/list projections and length;
\item Boolean \code{all}, \code{any}, and implication;
\item finite sum, product, maximum, and argmin/argmax;
\item first successful bounded search;
\item table lookup with a total default.
\end{itemize}

Each correctness theorem should be a \code{Prf PRAAx} object.  In particular,
\code{argmin} needs the invariant
\[
 \mathsf{bestIndex}<i+1\quad\wedge\quad
 \forall j<i+1\;v(\mathsf{bestIndex})\le v(j),
\]
where the bounded universal part is represented by a Boolean fold.  A
least-index tie rule makes the output deterministic.

\subsection{Finite query tables for names}

For a fixed evaluator symbol \code{a}, define formulas
\code{RegularA}, \code{EqualAB}, and \code{ApartAB}.  A theorem-specific query
planner produces a finite list of precision indices.  The object-language
\code{GoodTable} contains:

\begin{itemize}
\item every instantiated regularity inequality;
\item every domain and magnitude bound;
\item every finite relation between repeated samples;
\item all externally supplied certificate checks.
\end{itemize}

The proof of \code{RegularA -> GoodTable = 1} is separate from trace induction.
Universal elimination is applied at each generated index; a fold combines the
facts.  For adaptive queries, an induction on the trace stage derives the next
instance at the term computed by the current state.

\subsection{Rational and interval kernel}

A canonical object-language rational module should provide:

\begin{enumerate}
\item signed integers and positive-denominator rationals;
\item cross-multiplied equality and order;
\item exact addition, multiplication, absolute value, min/max, powers;
\item dyadics and a function finding a dyadic below a positive rational;
\item interval addition, multiplication with magnitude bounds, and guarded
      reciprocal;
\item finite polynomial evaluation and the Lipschitz majorant of
      \cref{lem:poly-lip}.
\end{enumerate}

The key target is not an abstract ordered-field hierarchy.  It is a small set
of concrete derivations used by the trace checker.  Ambient algebraic
hierarchies may help validate formulas, but the endpoint must use coded
natural-number operations.

\subsection{Robust branch and finite-net modules}

The overlapping branch theorem is a short embedded rational inequality proof.
Its implementation inputs codes \(a,e,q,d\), checks \(4e\le d\), compares
\(a\) with \(q+d/2\), and returns a branch bit plus one of two checked margin
inequalities.  The exact-gap theorem consumes an apartness certificate.

The finite-net module needs:
\begin{itemize}
\item grid generation and a finite coverage theorem for each standard domain;
\item bounded argmin correctness;
\item a Lipschitz transfer lemma;
\item nested-ball regularity and uniqueness.
\end{itemize}
Coverage is domain-specific.  Boxes and simplices should be implemented first;
sphere/projective charts can then reuse robust branch and normalization.

\subsection{Height and complete-list module}

Formalize \cref{thm:height-list} generically over fixed symbols for the point
predicate, height, and bounded enumerator.  The forward direction is a bounded
filter theorem.  The reverse direction is maximum-on-list correctness plus
list completeness.  This module provides a precise reusable endpoint for
Faltings, Siegel, integral-point, and finite-generation packets: they need only
supply the theorem-specific height or covering certificate.

\subsection{Certificate mode and public formulas}

A finite certificate language can be checked in ordinary \code{PRAAx} and
then translated.  If a smaller external proof format is desired, the guarded
public calculus in \code{NativePublic.lean} is relevant: equations are wrapped
by a gating, transport sources are protected, and the public translation is
interderivable with the native kernel on the intended language.  A future
proof compiler can emit public certificates while the trusted Lean theorem
proves their equivalence to \code{NPrf}.

This is an optimization, not the first milestone.  The initial implementation
should favor readable ordinary arithmetic derivations.

\subsection{Suggested module graph}

\begin{center}
\begin{tabular}{@{}l@{}}
\code{CFP/Codes.lean}\\
\quad$\downarrow$\\
\code{CFP/BooleanFold.lean}, \code{CFP/Rational.lean}\\
\quad$\downarrow$\\
\code{CFP/TraceReflection.lean}\\
\quad$\downarrow$\\
\code{CFP/RegularNames.lean}, \code{CFP/RobustBranch.lean}\\
\quad$\downarrow$\\
\code{CFP/FiniteNet.lean}, \code{CFP/PolynomialCertificate.lean},
\code{CFP/HeightSearch.lean}\\
\quad$\downarrow$\\
\text{domain packages and theorem-specific compilers}\\
\quad$\downarrow$\\
\code{pra_iff_native} and optional public-certificate export.
\end{tabular}
\end{center}

\subsection{Milestones and audits}

\begin{enumerate}
\item Compile the finite-trace semantic prototype and prove it without
      \code{sorry}.  This checks the abstract recurrence only.
\item Implement and compile the object-language Boolean-fold library.
\item Construct the first generic \code{PRA.traceReflection} proof.
\item Use it for the small recurrence and robust-branch examples in
      \cref{sec:examples}; inspect \code{\#print axioms}.
\item Formalize regular-name addition, multiplication, and reciprocal under
      apartness.  These are the first all-precision tests of the query-table
      discipline.
\item Formalize polynomial certificate reflection and finite-grid argmin.
\item Port one existing packet's rational core, preferably Euler--Rodrigues or
      finite probability, and replace theorem-specific induction plumbing by
      the generic modules.
\item Add height/list equivalence and port one explicit Diophantine family.
\item Only then optimize certificate size or target the public guarded
      calculus.
\end{enumerate}

Every milestone should run:
\begin{itemize}
\item a source scan for \code{sorry}, \code{admit}, and new line-leading
      \code{axiom};
\item \code{\#print axioms} on exported endpoints;
\item a clean build against the pinned Lean/Mathlib versions;
\item exact regression tests for independent semantic implementations;
\item a status ledger distinguishing semantic, architectural, and embedded
      completion.
\end{itemize}

\subsection{Artifacts supplied in this archive}

\code{lean/FiniteTraceSemantic.lean} gives an ordinary Lean prototype of the
trace invariant and complete-list-to-height argument.  It is not an embedded
proof.  \code{lean/IPRAObjectLanguageBlueprint.lean} imports the actual
\code{ConciseIPRA} interfaces, defines the proposed fixed-symbol record, and
states the completion path.  The nontrivial embedded theorem is deliberately
not claimed complete.  \code{lean/LEAN_PLAN.md} gives the dependency order and
target theorem names.  Lean was unavailable in the artifact-building
environment, so these files were not compiled.

\section{Worked finite examples}
\label{sec:examples}

The accompanying Python program is an exact reference implementation using
integers and \code{fractions.Fraction}.  It is not a proof in \(\ipra\), but
its examples mirror the finite recurrences to be embedded.

\subsection{A bounded trace}

Fix an input \(u\) and define
\[
 s_0=0,\qquad s_{i+1}=s_i+u+i.
\]
The invariant is
\[
 s_i=iu+\frac{i(i-1)}2.                         \tag{E1}\label{eq:trace-demo}
\]
The base case is immediate.  At successor,
\[
 s_{i+1}=iu+\frac{i(i-1)}2+u+i
 =(i+1)u+\frac{i(i+1)}2.
\]
Because \(i(i-1)\) is even, one may either use coded exact division by two or
avoid division by storing \(2s_i=2iu+i(i-1)\).  The latter is the cleaner
quantifier-free invariant.

For \(u=3\) and bound \(8\), the certified states are
\[
 0,3,7,12,18,25,33,42,52.
\]
The checker validates every local step and every instance of
\eqref{eq:trace-demo}.  The final value is \(52\).

\subsection{An overlapping decision}

Take
\[
 a=\frac9{16},\qquad e=\frac1{32},\qquad
 q=\frac12,\qquad d=\frac18.
\]
Here \(4e=d\), and
\(a=q+d/2\).  The upper branch of \cref{thm:overlap-branch} returns
\[
 x<q+d=\frac58.
\]
Indeed \(x\le a+e=19/32<20/32=5/8\).  Had the approximation been larger than
\(9/16\), the lower branch would certify \(x>1/2\).  The interval
\((1/2,5/8)\) is deliberately shared by both branches.

This tiny example is the complete control-flow core of much larger chart and
pivot arguments.

\subsection{Finite-grid approximation to \texorpdfstring{\(\sqrt2\)}{sqrt(2)}}

Let \(P(x)=x^2-2\) on \([1,2]\).  On this interval,
\cref{lem:poly-lip} gives Lipschitz constant \(4\).  Scan the grid with
denominator \(16\).  The least point minimizing \(\abs{P}\) is
\[
 x_* = \frac{23}{16},
 \qquad
 \abs{P(x_*)}=\frac{17}{256}.
\]
A nearest grid point to \(\sqrt2\) is within \(1/32\), so the generic net
bound gives residual at most \(4/32=1/8\); the exact residual above is smaller.
The scan needs no square-root function.  It uses only integer squares and
rational comparison.

To obtain an exact regular name, at precision \(k\) compute
\[
 n_k=\floor{2^{k}\sqrt2}
     =\mathsf{isqrt}(2\cdot2^{2k}),
\]
checked by
\[
 n_k^2\le2\cdot2^{2k}<(n_k+1)^2.
\]
The interval
\([n_k/2^k,(n_k+1)/2^k]\) has width \(2^{-k}\), contains \(\sqrt2\), and is
nested up to a harmless index shift.  Integer square root is a bounded search
with a quantifier-free invariant.  The archive's regression test checks the
resulting midpoint regularity through precision twelve; the general proof is
the nested-ball theorem.

\subsection{Polynomial identity certificates}

Consider
\[
 (x+y)^2=x^2+2xy+y^2.
\]
A certificate producer expands the left side, sorts monomials
\((2,0),(1,1),(0,2)\), and returns coefficients \((1,2,1)\).  The right side
already has that normal form.  The checker verifies each list transformation.
Soundness is proved once by induction on the certificate list; the theorem for
this identity is then one accepted finite certificate plus universal
introduction of \(x,y\).

For a rational identity, the checker also records a positive common
denominator.  The final statement is the corresponding integer polynomial
identity after multiplication.  This is the recommended path for matrix,
trigonometric, and algebraic identities in the packets.

\subsection{Height bound and list}

Let \(P(a,x)\) be \(x^2\le a\), with height \(h(x)=x\).  For \(a=20\), the
height certificate \(H=4\) produces the complete list
\([0,1,2,3,4]\).  Conversely the maximum of the list is \(4\).  The example is
trivial, but its proof structure is identical to bounded rational-point search:
a theorem-specific height implication feeds a generic finite filter.

\subsection{Finite probability}

Let a law have rational values \(x_j=a_j/b_j\) and natural weights \(w_j\),
with \(W=\sum_jw_j>0\).  Its expectation is the exact rational
\[
 \mathbb E X=\frac{1}{W}\sum_jw_jx_j.
\]
For two laws, enumerate pairs by mixed radix.  The identity
\[
 \sum_{i,j}w_i v_j f(x_i)g(y_j)
 =\left(\sum_iw_i f(x_i)\right)
  \left(\sum_jv_j g(y_j)\right)
\]
is a finite polynomial certificate.  It proves independence algebra, mean
additivity, and variance addition for centered variables.  There is no measure
space and no appeal to Fubini; all sums are finite and their indices have
explicit bounds.

\subsection{Reproduction}

From the archive root, run
\begin{verbatim}
cd code
python3 -m unittest -v test_certificate_kernel.py
python3 certificate_kernel.py ../examples/framework_certificate.json
\end{verbatim}
The delivered run passes nine tests.  The JSON certificate records the trace,
branch, finite-grid, and regularity examples.  These tests provide independent
finite arithmetic checks; they do not replace the planned \code{Prf PRAAx}
formalization.

\section{Further research questions}
\label{sec:research}

\begin{enumerate}
\item \textbf{Generic object-language trace reflection.}
Can the proposed \code{PRTraceSystem} record be implemented with sufficiently
clean arity and substitution lemmas that theorem-specific modules only supply
base, step, and final derivations?

\item \textbf{Verified certificate compiler.}
Can a small external compiler translate polynomial normalizations, rational
error calculations, and finite loop annotations into public guarded
certificates whose checker is proved equivalent to \code{NPrf}?

\item \textbf{Canonical code library.}
Which existing ProveIt encodings of Booleans, pairs, lists, integers, and
rationals should become canonical, so that proof packets stop rebuilding
incompatible arithmetic kernels?

\item \textbf{Automatic precision synthesis.}
For expression trees with sums, products, guarded inverses, roots, and finite
folds, can Lean synthesize a monotone primitive-recursive precision transformer
and emit the corresponding error proof?

\item \textbf{Adaptive query extraction.}
Can a theorem-specific semantic program be instrumented to emit its finite
query trace, after which a generic object-language checker proves that all
queried regularity instances suffice?

\item \textbf{Constructive root clusters.}
Develop a complete isolating-cluster library for integer polynomials, including
multiplicity, subdivision certificates, and exact coefficient reconstruction,
without imposing a discontinuous global root order.

\item \textbf{Projective linear algebra.}
Develop projective kernel, eigenspace projector, and singular-value residual
interfaces that give unconditional approximate outputs and exact outputs under
quantitative gaps.  Compare their proof size with classical normalized-vector
interfaces.

\item \textbf{Effective compact domains.}
Formalize certified nets for rational boxes, simplices, spheres, projective
spaces, compact matrix groups, and finite probability simplices.  Determine a
minimal common interface for coverage and refinement.

\item \textbf{Analytic certificate formats.}
What finite certificate best represents zero-free regions, contour bounds,
exceptional zeros, or saddle-point estimates so that downstream analytic
number theory becomes a small verifier theorem?

\item \textbf{Height-certificate taxonomy.}
Classify Diophantine methods---Runge, descent, Baker bounds, Chabauty,
Mordell--Weil sieve---by the finite certificate they can produce and the
checker needed in \(\ipra\).

\item \textbf{Witness extraction bounds.}
For restricted formula classes proved in ordinary or native \(\ipra\), obtain
explicit primitive-recursive bounds on extracted witnesses and compare them
with direct \(\CFP\) producers.

\item \textbf{Obstruction automation.}
Create reusable reductions from exact sign, exact root ordering, exact chart
selection, or negative-only apartness to LLPO, LPO, or Markov principles.  A
packet could then test a proposed classical formulation before investing in a
large positive proof.

\item \textbf{Proof-size engineering.}
Measure whether ordinary \code{PRAAx} certificates, native formulas, or guarded
public certificates give the smallest checker and proof terms for large
finite mathematics.

\item \textbf{Repository-wide status ledger.}
Can every proof packet expose a machine-readable record following
\code{metadata/framework_schema.json}, allowing automated reports of which
fields---producer, verifier, embedded proof, native export, axiom audit---are
complete?

\item \textbf{A substantial end-to-end pilot.}
A strong first pilot would combine regular names, robust pivots, polynomial
certificates, and finite search: for example, the rational and separated-real
Euler--Rodrigues theorems or a finite-law Berry--Esseen core.  Completing one
such packet with actual \code{Prf PRAAx} endpoints would test the entire
architecture.
\end{enumerate}

\appendix
\section{Object-language normal forms and packet checklist}
\label{app:normal-forms}

\subsection{Finite exact endpoint}

For input code \(p\), algorithm \(A\), and Boolean checker \(V\):
\[
 \forall p\quad \mathsf{Dom}(p)=1\to V(p,A(p))=1.
\]
Required proof objects:
\begin{enumerate}
\item \code{PR} symbols for \(\mathsf{Dom},A,V\);
\item defining-equation lemmas;
\item QF loop invariant;
\item checker soundness implication;
\item closed \code{Prf PRAAx} theorem and native translation.
\end{enumerate}

\subsection{Regular-name endpoint}

For fixed evaluator symbols \(a,b\):
\[
 \forall p\quad
 \mathsf{Reg}_a(p)\wedge\mathsf{Dom}(p)
 \to \forall k\;\mathsf{Close}(b(p,k),\mathsf{Target}(p),k).
\]
Here \(\mathsf{Target}\) must itself be expressed through finite observations,
or the conclusion should be equality of names.  The proof at free \(k\) must
list every queried input precision.

\subsection{Approximate existential endpoint}

Instead of \(\exists y\,P(p,y)\) in a higher-type completion, use
\[
 \forall k\;\exists y<B(p,k)\quad
 \mathsf{Residual}(p,y)\le2^{-k}.
\]
A stronger computational form provides a fixed producer \(Y(p,k)\):
\[
 \forall p,k\quad
 \mathsf{Dom}(p)\to
 \mathsf{Residual}(p,Y(p,k))\le2^{-k}.
\]
Coherence of \(Y(p,k)\) is a separate theorem; without it no exact name is
claimed.

\subsection{Exact existential from nested certificates}

Input or construct \(B(p,k)\), a nested family of finite regions, with
checked existence and shrinking radius.  Define the output name from their
centers.  Prove regularity, residual convergence, and uniqueness if supplied.
The object-language theorem quantifies only over natural codes and precision.

\subsection{Certificate-relative endpoint}

\[
 \forall p,c\quad
 \mathsf{CertOK}(p,c)=1\to
 \mathsf{Conclusion}(p,\mathsf{decode}(c)).
\]
The packet must state whether a certificate producer is:
\begin{itemize}
\item proved uniformly primitive recursive;
\item proved for a restricted input class;
\item executable but not embedded;
\item assumed as explicit input;
\item unknown or obstructed.
\end{itemize}

\subsection{Status vocabulary}

\begin{longtable}{@{}p{0.24\textwidth}p{0.68\textwidth}@{}}
\toprule
Status & Meaning \\
\midrule
\endhead
Specified & Exact codes, formulas, and proof obligations are written.\\
Finite core proved & Finite semantic algorithms and identities are checked,
but no full embedded endpoint exists.\\
Embedded proof complete & A \code{Prf PRAAx} endpoint compiles and has an
axiom audit.\\
Conditional & Verifier and consequence are proved relative to a finite
certificate or modulus whose producer is separate.\\
Obstructed & A stronger uniform statement implies an excluded constructive
principle or noncomputable function; a maximal positive replacement is given.\\
\bottomrule
\end{longtable}

\subsection{Review checklist}

Before accepting a new packet, answer:
\begin{enumerate}
\item What is the exact first-order formula?
\item Which function symbols are fixed, and which data are numerical inputs?
\item Is any purported universal evaluator hidden in the notation?
\item What is the requested precision and its dependency graph?
\item Which denominators, pivots, or charts require positive margins?
\item What is the primitive-recursive loop/search bound?
\item What is the quantifier-free invariant, after Booleanization?
\item Are all compactness/existence uses replaced by nets, bounds, or nested
      certificates?
\item Is an ineffective constant being used as though it were produced?
\item Is the Lean result semantic, embedded, or merely a target proposition?
\item Has the final declaration been compiled and audited for axioms?
\item If the literal classical statement fails constructively, is the maximal
      positive replacement and obstruction both recorded?
\end{enumerate}

\section{Repository source audit}
\label{app:sources}

\subsection{Snapshot}

Repository: \url{https://github.com/VladimirReshetnikov/ProveIt}

Inspected main commit:
\begin{center}
\small\ttfamily
7d492fbfc5fee3b2b9dec0b2623b0b82effe05f2.
\end{center}
The inspection was read-only through the GitHub connector.  A network clone
was unavailable in the artifact runtime, so the report relies on the fetched
files listed below rather than claiming an exhaustive local build.

\subsection{Foundation files}

\begin{longtable}{@{}p{0.49\textwidth}p{0.43\textwidth}@{}}
\toprule
Path & Use in this report \\
\midrule
\endhead
\code{Logic/PrimitiveRecursiveArithmetic/ConciseIPRA/README.md} & Native term
and formula grammar, rule overview, and design intent.\\
\code{.../ConciseIPRA/Lean/ConciseIPRA/PR.lean} & \code{PR}, \code{PTm},
\code{PRAAx}, \code{PRA.qfInd}, primitive-recursion equations.\\
\code{.../ConciseIPRA/Lean/ConciseIPRA/Native.lean} & \code{NFm}, \code{NPrf},
substitution and soundness.\\
\code{.../ConciseIPRA/Lean/ConciseIPRA/BiInterpretation.lean} & Faithful
translation between equality theory and iPRA.\\
\code{.../ConciseIPRA/Lean/ConciseIPRA/NativeBiInterpretation.lean} &
\code{pra_iff_native}, \code{native_iff_pra}, round trips.\\
\code{.../ConciseIPRA/Lean/ConciseIPRA/NativePublic.lean} & Guarded public
certificate language and substitution-stable proof restrictions.\\
\bottomrule
\end{longtable}

\subsection{Proof-packet index and representative audits}

\begin{longtable}{@{}p{0.51\textwidth}p{0.41\textwidth}@{}}
\toprule
Path & Finding used \\
\midrule
\endhead
\code{Logic/PrimitiveRecursiveArithmetic/docs/proofs/README.md} & Forty-eight
packets; placement did not rerun or compile them; area-by-area status and
explicit partial results.\\
\code{.../geometry/euler_rodrigues_ipra/README.md} & LLPO obstruction to
unrestricted exact unit-axis choice; rational, separated, approximate, and
quaternion positive results.\\
\code{.../geometry/euler_rodrigues_ipra/formalization_and_audit.tex} & Fix
precision before induction; finite \code{GoodTable}; polynomial residual and
precision-compilation discipline.\\
\code{.../diophantine-number-theory/dirichlet_kronecker_ipra/README.txt} &
LLPO/LPO/Markov obstructions and strict/separated positive approximation
forms.\\
\code{.../analytic-number-theory/zeta_ipra_proof_package/paper/sections/01_foundations.tex} &
Regular names, fixed evaluator schemas, no universal PR evaluator, finite
mesh/series discipline.\\
\code{.../analytic-number-theory/prime_distribution_ipra/paper/sections/01_scope.tex} &
Explicit bounds for positive results; incomplete Chebotarev,
Siegel--Walfisz, and Bombieri--Vinogradov analytic inputs.\\
\code{.../diophantine-number-theory/faltings_ipra/README.md} & General
Faltings not proved; height bound versus list; cardinality does not give a
stopping rule; ambient Lean versus embedded proof.\\
\code{.../diophantine-number-theory/siegel_ipra/paper/01_scope.tex} & General
Siegel not proved; explicit Runge certificate families; missing general size
bound.\\
\code{.../diophantine-number-theory/mordell_weil_ipra/README.txt} & Generic
finite descent relative to actual covering representatives; no uniform
arbitrary-curve generator algorithm.\\
\code{.../probability/central_limit_ipra_2026-11-15/paper/central_limit_ipra.tex} &
Finite rational laws, effective completion with moduli, rational thresholds,
and a Lean blueprint rather than a completed embedded proof.\\
\bottomrule
\end{longtable}

\subsection{Artifact verification}

The exact Python companion in this archive passed nine unit tests and verified
the supplied JSON certificate.  The PDF was compiled from the supplied LaTeX,
rendered to images, and visually inspected.  The Lean files were not compiled
because no Lean executable was available.  Their status is therefore exactly
``semantic prototype and object-language blueprint.''


\end{document}
